\documentclass[11pt,a4paper]{article}
\usepackage[T1]{fontenc}
\usepackage[utf8]{inputenc}
\usepackage[margin=25mm]{geometry}
\usepackage{lmodern}
\usepackage{microtype}
\usepackage{xcolor}
\newif\ifmarkrevisions
\markrevisionsfalse
\ifmarkrevisions
  \newcommand{\rev}[1]{{\color{blue}#1}}
  \newenvironment{revision}{\color{blue}}{}
\else
  \newcommand{\rev}[1]{{\color{black}#1}}
  \newenvironment{revision}{\color{black}}{}
\fi
\newcommand{\ess}{\mathop{\mathrm{ess}}}
\newcommand{\doi}[1]{\href{https://doi.org/#1}{\nolinkurl{doi:#1}}}
\newcommand{\locmark}[1]{\label{#1}\linelabel{line:#1}}
\setlength{\emergencystretch}{2em}

\usepackage{amsfonts}
\usepackage{amssymb,amsthm,amsmath}
\usepackage{enumitem}
\usepackage{multicol}


\newcommand{\id}{\mathop{\mathrm{id}}}
\newcommand{\stab}{\mathop{\mathrm{stab}}\nolimits}
\newcommand{\supp}{\mathop{\mathrm{supp}}}
\newcommand{\type}{\mathop{\mathrm{type}}}
\newcommand{\Sym}{\mathop{\mathrm{Sym}}}

\newtheorem{prop}{Proposition}
\newtheorem{deft}[prop]{Definition}

\usepackage{orcidlink}
\usepackage{hyperref}
\hypersetup{hidelinks}
\pdfstringdefDisableCommands{\def\rev#1{#1}}
\usepackage[mathlines]{lineno}
\renewcommand\linenumberfont{\normalfont\tiny\color{black}}
\modulolinenumbers[5]


\begin{document}
\title{\rev{Irreducible Local Rules and Geometric Equivalence Classes of One-Dimensional Cellular Automata}}

\author{
    Martin Schaller\,\orcidlink{0009-0000-1939-8475}\thanks{martin.roman.schaller@gmail.com} \\
    \small{Vienna, Austria}
    \and
    Karl Svozil\,\orcidlink{0000-0001-6554-2802}\thanks{karl.svozil@tuwien.ac.at, \url{http://tph.tuwien.ac.at/~svozil}} \\
    \small{Institute for Theoretical Physics, TU Wien,} \\
    \small{Wiedner Hauptstrasse 8-10/136, 1040 Vienna, Austria}
}

\date{27 October 2026}
\maketitle
\linenumbers

\begin{revision}
\begin{abstract}
We count geometric equivalence classes of one-dimensional cellular automata
by combining irreducible local rules with the geometry of their essential
neighbourhoods. Using the standard uniqueness of minimal memory sets, we
derive recurrences for the stabilizer types of irreducible binary rules and
enumerate global maps under state permutations, reflection, and
postcomposition by shifts. Asymmetric essential neighbourhoods require an
additional orientation term in the counting formula. We define rational
neighbourhood scaling intrinsically as a partial operation, prove that the
resulting similarity relation is an equivalence relation, and describe integer
scaling as conjugacy to a product of independent copies. For binary rules on
three- and four-site intervals, the congruence counts are respectively
85 and 16\,656; allowing scaling gives 81 and 16\,648. These counts are also
verified by exhaustive enumeration. On the finite set of elementary cellular
automata, the join of similarity and topological conjugacy has 76 classes.
\end{abstract}

\section{Introduction}
\locmark{loc:introduction}
One-dimensional cellular automata (CAs) update a bi-infinite configuration
over a finite alphabet by applying the same local rule at every lattice site.
The global map determines the dynamics, but its local description is not
unique: a neighbourhood may contain sites that the rule does not use.
For classification, this redundancy must be separated from changes in the
position, orientation, and scale of the sites that are actually used.

The Curtis--Hedlund--Lyndon theorem characterizes CA global maps as the
continuous maps that commute with the shift
\cite{DBLP:journals/mst/Hedlund69,Lind_Marcus_1995}.
Every CA also has a unique minimal memory set
\cite[Section~1.5]{CeccheriniCoornaert2010}. This set consists exactly of
the essential coordinates of a local representation
\cite[Proposition~1]{CastilloVeliz2024}.
We call a local rule irreducible when it depends on every input coordinate.
Minimal-memory uniqueness then gives a canonical pair consisting of an
essential neighbourhood and an irreducible rule. We use this established
result as a foundation for enumeration, and include a short proof for
completeness.

\locmark{loc:prior-work}
Reflection and state relabeling give the familiar finite-group classification
of local rules. For elementary cellular automata (ECAs), the 256 binary
three-input rules form 88 classes
\cite{wolfram-1986,Li1990TheSO}.
Cattaneo et al.\ \cite{CATTANEO19971593} studied transformations of the
one-dimensional rule space. Our earlier work
\cite{SCHALLER2025100298} counted symmetry orbits for two- and three-state
rules and refined them by stabilizer type. Its Section~3.7 already defined
irreducibility and discussed shift-equivalence using global maps, while
leaving the integration of shifts into the orbit-counting framework outside
its scope. The present paper develops that integration and accounts for
asymmetric essential supports.

Shifting local rules and neighbourhoods has also been studied by
Acerbi et al.\ \cite{10.1007/978-3-540-73001-9_1},
Ruivo et al.\ \cite{RUIVO2018280}, and
Garc{\'i}a-Morales \cite{GARCIAMORALES2013276}.
More general algebraic settings are considered in
\cite{CASTILLORAMIREZ2020104533,Castillo-Ramirez03072023}.
Here a shift of the essential neighbourhood acts by
\emph{postcomposition}, $\Phi\mapsto\sigma^a\Phi$.
It is not conjugation by $\sigma^a$: shift commutation gives
$\sigma^a\Phi\sigma^{-a}=\Phi$.
Thus our geometric equivalences need not preserve the dynamical system up
to topological conjugacy.

\locmark{loc:contributions}
We distinguish the established ingredients from the contributions developed
here as follows.
\begin{enumerate}
\item Minimal-memory uniqueness and the elementary inclusion--exclusion
count for irreducible rules are recalled in Section~\ref{sec:irr}.
The unrestricted local-rule orbit counts and stabilizer counts are taken
from \cite{SCHALLER2025100298}.
\item Section~\ref{sec:loc} derives recurrences that extract the
\emph{irreducible} stabilizer counts from those unrestricted counts
(Proposition~\ref{prop:recurrences}).
\item Section~\ref{sec:glob} combines these counts with neighbourhood
geometry. Proposition~\ref{prop:global-eq} describes the orientation
splitting for asymmetric supports, and
Proposition~\ref{prop:global-shift} gives the resulting congruence count.
An independent exhaustive enumeration verifies the small-arity results.
\item Propositions~\ref{prop:scaling-domain} and
\ref{prop:similarity} formulate rational scaling on its intrinsic domain
and establish similarity as an equivalence relation.
Proposition~\ref{prop:global-scaling} gives the corresponding count.
For integer scaling, Proposition~\ref{prop:product} identifies the precise
product dynamics.
\end{enumerate}

These results concern geometric rule classification. They complement
dynamical classifications; see \cite{VISPOEL2022133074} for a survey.
In particular, Epperlein's classification of ECAs up to topological
conjugacy has 83 classes
\cite{epperlein2015classification,epperlein2017}.
We compare that partition with our 81 similarity classes and compute their
join on the ECA set. Section~\ref{sec:extended} explains the distinctions
between conjugacy, moving frames, product dynamics, and factor maps.
\end{revision}

\section{Cellular Automata}
\label{sec:CA}

If $A$ is any set, we denote the size (or cardinality) of $A$ by $|A|$.
We write $\mathcal{P}(A)$ for the power set of $A$, the set of all subsets of $A$.
\rev{Throughout, the alphabet $\mathcal{A}$ is finite and has $k\geq2$ symbols.}
A word $w=x_1x_2\ldots x_{m}$ over the alphabet $\mathcal{A}$ is a finite sequence of symbols
from $\mathcal{A}$ juxtaposed, and we write $|w|$ to \rev{denote the length} of the sequence,
that is $|x_1x_2\ldots x_{m}| = m$
(the notation $|.|$ denotes both the size of a set and the length of a word).
We write $\mathcal{A}^m$ for the set of all words of length $m$ over the alphabet $\mathcal{A}$.
We denote the integer interval $\{p, p+1, \ldots, q\}$ by $[p,q]$.
If $q < p$, then $[p,q]$ is the empty set.
A configuration $x$ is a bi-infinite sequence over
the alphabet $\mathcal{A}$, formally defined as a map $\mathbb{Z} \to \mathcal{A}$.
The \rev{$i$th} element, $i \in \mathbb{Z}$, of a configuration $x$ is denoted by $x_i$.


\begin{deft}
\label{def:CA}
A one-dimensional CA $A$ is a triple $(\mathcal{A}, N, f)$, where \\
\rev{$\mathcal{A}$ is a finite set of at least two states;} \\
$N$ is the neighbourhood, a finite set of integers of size $n$; \\
$f$ is the local rule, a function from $\mathcal{A}^n$ to $\mathcal{A}$. \\
Let $N=\{j_1,j_2, \ldots,j_{n}\}$ such that $j_1 < j_2 < \cdots < j_{n}$.
The neighbourhood $N$ and the local rule $f$ induce the global map
$\Phi^A : \mathcal{A}^\mathbb{Z} \rightarrow \mathcal{A}^\mathbb{Z}$ on the set of configurations, defined by
$\Phi^N_f(x)_{i} = f(x_{i+j_1}x_{i+j_{2}}\ldots x_{i+j_{n}})$.
\end{deft}
We will also use the notation $\Phi^N_f$ for the global map, if the state set $\mathcal{A}$ is given.
If the CA is initialised with the configuration $x$, the CA computes in one step
the configuration $\Phi^N_f(x)$.
In case of $\mathcal{A} = \{0,1\}$, the CA, the local rule, and the global map, are all called binary.
The shift map (or operator) on $\mathcal{A}^\mathbb{Z}$ maps a configuration $x$ to the configuration $\sigma(x)$
whose $i$th coordinate is $\sigma(x)_i = x_{i+1}$.
It follows from the definition of the CA, that the global map commutes with the shift operator, that is
$\Phi^N_f(\sigma(x)) = \sigma(\Phi^N_f(x))$.

\begin{revision}
For distinct configurations $x,y\in\mathcal{A}^{\mathbb Z}$, put
\[
d(x,y)=2^{-\min\{|i|:x_i\ne y_i\}},\qquad d(x,x)=0.
\]
This metric induces the product topology on $\mathcal{A}^{\mathbb Z}$.
The alphabet is equipped with the discrete topology.
\end{revision}

The following proposition, see \cite{DBLP:journals/mst/Hedlund69}, shows that CAs can also be defined by topological means.
\begin{prop}[\rev{Curtis--Hedlund--Lyndon} Theorem]
A map of $\mathcal{A}^\mathbb{Z}$ into $\mathcal{A}^\mathbb{Z}$ is the global map of a CA if and only if
the map is continuous and commutes with the shift operator.
\end{prop}
\rev{A presentation of this theorem and symbolic dynamics is given
in~\cite{Lind_Marcus_1995}.}
\begin{revision}
\locmark{loc:dynamical-definitions}
For continuous maps $\Phi:X\to X$ and $\Psi:Y\to Y$, a continuous
map $h:X\to Y$ satisfying $h\circ\Phi=\Psi\circ h$ is a
homomorphism of dynamical systems \cite[p.~185]{Lind_Marcus_1995}.
Terminology for semiconjugacy varies; here we use \emph{semiconjugacy},
or \emph{factor map}, for a surjective homomorphism. A topological
conjugacy is a homeomorphism satisfying the same intertwining identity.
A topological conjugacy between CA dynamical systems is not required
to commute with the spatial shift.
\end{revision}
We call the set $L_\mathcal{A}(n) = \{f \, | \, f : \mathcal{A}^n \to \mathcal{A} \}$ the local rule space of
arity $n$ over the alphabet $\mathcal{A}$.
If we set $k = |\mathcal{A}|$, the size of $L_{\mathcal{A}}(n)$ is $k^{k^n}$.
The set $G_\mathcal{A}(N) = \{ \Phi^N_f \, | \, f \in L_\mathcal{A}(|N|) \}$ is called the global rule
space on the neighbourhood $N$ over the alphabet $\mathcal{A}$.
Wolfram~\cite{wolfram-1986} introduced a commonly used coding schema to enumerate the local rule space $L_{\{0,1\}}(n)$.
Let $f \in L_{\{0,1\}}(n)$ be a binary local rule.
The Wolfram code $c(f)$ of $f$ is the nonnegative integer
\[
c(f) = \sum_{a_1\ldots a_{n} \in \{0,1\}^n} f(a_1\ldots a_{n}) \times 2^{\sum_{i=1}^{n} a_i 2^{n-i}}.
\]
This summation uses the left-to-right most significant bit (MSB) convention.
If $f: \{ 0, 1 \}^n \to \{ 0, 1 \} $ is a local rule,
we define $W^n_{c(f)} = f$.
We denote the constant functions $f() = 0$ and $g() = 1$, both of arity 0, by
$W^0_0$ and $W^0_1$, respectively.
If this coding schema is used in combination with a global map, the arity of $f$ can be derived
from the neighbourhood, for instance, we can write
$\Phi^{[-1,1]}_{W_{110}}$ instead of $\Phi^{[-1,1]}_{W^3_{110}}$.

A global map of a CA can have different representations in terms of
local rules and neighbourhoods.
The binary identity map $\id: \{0,1\}^\mathbb{Z} \to \{0,1\}^\mathbb{Z}$,
$\id(x) = x$ serves as a trivial example.
For each neighbourhood $N$ that contains the integer 0,
the map $\id$ has a different representation, so that for instance
\[
\id = \Phi^{\{0\}}_{W_2} = \Phi^{[0,1]}_{W_{12}} = \rev{\Phi^{[-1,1]}_{W_{204}}}.
\]
\rev{More generally, if $N_2 \subseteq N_1$} then $G_\mathcal{A}(N_2) \subseteq G_\mathcal{A}(N_1)$.
\rev{In the next section} we will see that each global map of a CA
has a unique canonical representation.

\section{Irreducibility}
\label{sec:irr}
\begin{revision}
\locmark{loc:minimal}
The uniqueness of a minimal memory set is standard
\cite[Section~1.5, Propositions~1.5.1--1.5.2]{CeccheriniCoornaert2010};
see also \cite[Proposition~1]{CastilloVeliz2024} for its characterization by
essential coordinates. Propositions~\ref{prop:uniqueness} and
\ref{prop:eq-cont} recall this fact in the notation used here; they are
background results, not an inverse or extension of the
Curtis--Hedlund--Lyndon theorem.
\end{revision}
A local rule $f : \mathcal{A}^n \to \mathcal{A}$ is said to be dependent on index $j$, $1 \le j \le n$,
if there is a word $a_1\ldots a_{j-1}a_{j+1}\ldots a_{n} \in \mathcal{A}^{n-1}$ and
symbols $b, c \in \mathcal{A}$ such that
$f(a_1\ldots a_{j-1}ba_{j+1}\ldots a_{n}) \ne f(a_1\ldots a_{j-1}ca_{j+1}\ldots a_{n})$,
otherwise $f$ is said to be independent of the index $j$.
The local rule $f: \mathcal{A}^n \to \mathcal{A}$ is said to be irreducible
if $f$ depends on all its indices, otherwise
$f$ is said to be reducible.
A CA $(\mathcal{A}, N, f)$ is said to be irreducible or reducible, if its local rule $f$ is
irreducible or reducible, respectively.
Define
 \[
\supp(f) = \{ j \, | \, \mbox{$f$ depends on coordinate $j$} \}.
\]
Then the local rule $f: \mathcal{A}^n \to \mathcal{A}$ is irreducible if and only if
$|\supp(f)| = n$.

A map $h: \mathcal{A}^n \to \mathcal{A}^m$, $m \le n$, is called a projection if
there are integers $p_1, \ldots, p_{m}$ with $1 \leq p_1 < \ldots < p_{m} \leq n$ such that
$h(a_1\ldots a_{n}) = a_{p_1}\ldots a_{p_{m}}$ for all words $a_1\ldots a_{n} \in \mathcal{A}^n$.
The set $\{p_1, \ldots, p_{m}\}$ is called the index set of the projection.
The map $H$ associates a projection with its index set, that is $H(h) = \{p_1, \ldots, p_{m}\}$.
The number of distinct projections $h : \mathcal{A}^n \to \mathcal{A}^m$ is given by
the binomial coefficient $\binom{n}{m}$.
If $f: \mathcal{A}^n \to \mathcal{A}$ is a local rule and $|\supp(f)| = m$,
then there is an irreducible local rule $g: \mathcal{A}^m \to \mathcal{A}$ and a projection
$h : \mathcal{A}^n \to \mathcal{A}^m$ such
that $f = g \circ h$.
This representation is unique.
This means that if $f = g^\prime \circ h^\prime$, $g^\prime$ is irreducible, and $h^\prime$ is a projection,
then $g = g^\prime$ and $h = h^\prime$.
Define $\overline{L}_\mathcal{A}(n) = \{f \in L_\mathcal{A}(n) \, | \, \mbox{$f$ is irreducible}  \}$ and
$\overline{G}_\mathcal{A}(N) = \{ \Phi_f^N \ | \ f \in \overline{L}_\mathcal{A}(|N|) \}$.
\begin{prop}
\label{prop:uniqueness}
If $A=(\mathcal{A}, N_1, f_1)$ and $B=(\mathcal{A}, N_2, f_2)$ are two irreducible CAs,
then $\Phi^A = \Phi^B$ if and only if $A=B$, that is if $N_1 = N_2$ and $f_1 = f_2$.
\end{prop}
\begin{proof}
If $N_1 = N_2$, then $\Phi^A = \Phi^B$ implies $f_1 = f_2$.
We therefore assume that $N_1 \neq N_2$.
Consequently, $N_1 \setminus N_2$ is nonempty, or $N_2 \setminus N_1$ is nonempty.
Without loss of generality, suppose that $N_1 \setminus N_2$ is nonempty.
Let $N_1 = \{j_1, \ldots, j_{n} \}$ and suppose $j_r \in N_1 \setminus N_2$.
The irreducibility of $f_1$ implies that there are words
$w_1 = a_1 \ldots a_{r-1}ba_{r+1} \ldots a_{n}$ and
$w_2 =a_1 \ldots a_{r-1}ca_{r+1} \ldots a_{n}$ in $\mathcal{A}^n$
such that $f_1(w_1) \neq f_1(w_2)$.
Construct two configurations $x$ and $y$ as follows.
Fix a symbol $d \in \mathcal{A}$.
If $i \not\in N_1$ set $x_i = y_i = d$.
If $j_i \in N_1$ and $j_i \neq j_r$ set $x_{j_i} = y_{j_i} = a_i$.
Set $x_{j_r} = b$ and $y_{j_r} = c$.
From the construction of $x$ and $y$ follows that $\Phi^A(x)_0 = f_1(w_1) \neq f_1(w_2) = \Phi^A(y)_0$ and
$\Phi^B(x)_0 = \Phi^B(y)_0$.
Thus $\Phi^A(x) \neq \Phi^B(x)$ or $\Phi^A(y) \neq \Phi^B(y)$.
\end{proof}
\rev{Proposition~\ref{prop:uniqueness}, the irreducible factorization above, and the Curtis--Hedlund--Lyndon theorem imply the next proposition.}
\begin{prop}
\label{prop:eq-cont}
If  $\Psi: \mathcal{A}^\mathbb{Z} \to \mathcal{A}^\mathbb{Z}$ is a continuous map that commutes with the shift operator $\sigma$,
then there is exactly one irreducible CA $A$ such that $\Psi = \Phi^A$.
\end{prop}
\begin{revision}
\locmark{loc:essential}
Let $\mathcal C_{\mathcal A}$ be the set of continuous maps on
$\mathcal A^{\mathbb Z}$ that commute with the shift, equivalently
the set of all CA global maps over $\mathcal A$.
By Proposition~\ref{prop:eq-cont}, it is in bijection with the irreducible
CAs over that alphabet.
For $\Phi\in\mathcal C_{\mathcal A}$, write $E(\Phi)$ for its unique
essential neighbourhood and $f_\Phi$ for its irreducible local rule.
If $\Phi=\Phi_f^N$ and $N=\{j_1<\cdots<j_n\}$, then
\[
E(\Phi)=\{j_i:i\in\supp(f)\}.
\]

\end{revision}

\begin{prop}\label{prop:inclusion}
If $k=|\mathcal{A}|$, the size of $\overline{L}_\mathcal{A}(n)$ is given by $\sum_{i=0}^n (-1)^{n-i} \binom{n}{i} k^{k^i}$.
\end{prop}
\begin{proof}
Fix the alphabet $\mathcal{A}$ and put $k = |\mathcal{A}|$.
Set $a_n = |L_\mathcal{A}(n)| = k^{k^n}$ and $b_n = |\overline{L}_\mathcal{A}(n)|$.
Let $j \leq n$ and $f \in \overline{L}_\mathcal{A}(j)$.
Then $f$ contributes $\binom{n}{j}$ times to $a_n$.
We obtain the recurrence relation
\[
a_n = \sum_{j=0}^n \binom{n}{j} b_j.
\]
We plug the asserted equation
\[
b_j = \sum_{i=0}^j (-1)^{j-i} \binom{j}{i} a_i
\] into this relation and get
\begin{eqnarray*}
a_n = \sum_{j=0}^n \binom{n}{j} \sum_{i=0}^j (-1)^{j-i} \binom{j}{i} a_i
= \sum_{i=0}^n a_i \sum_{j=i}^n (-1)^{j-i} \binom{n}{j} \binom{j}{i}.
\end{eqnarray*}
The inner sum is known to satisfy the identity $\sum_{j=i}^n (-1)^{j-i} \binom{n}{j} \binom{j}{i} = \delta_{in}$, where $\delta_{in}$ is the Kronecker delta (equal to 1 if $i=n$, and 0 otherwise) \cite[p.~169]{Graham94}.
Thus the right hand side collapses to $a_n$, verifying the formula.
\end{proof}

\begin{table}
\begin{center}
\caption{Number of binary local rules.}
\label{tab:irr}
\begin{tabular}{|l|c|c|c|c|c|c|} \hline
$n$       & 0 & 1 & 2  & 3   & 4       & 5 \\ \hline
$|\overline{L}_2(n)|$   & 2 & 2 & 10 & 218 & 64\,594 & 4\,294\,642\,034\\ \hline
$|L_2(n)|$              & 2 & 4 & 16 & 256 & 65\,536 & 4\,294\,967\,296\\ \hline
\end{tabular}
\end{center}
\end{table}

The cardinalities of the binary local rule spaces for a small number of arguments are
depicted in Table~\ref{tab:irr}.
Let $H_\mathcal{A}(n, m)$ denote the set of projections $h : \mathcal{A}^n \to \mathcal{A}^m$.
The definition of an irreducible function implies the first part of the following proposition,
\rev{using $\mathcal A^0=\{\varepsilon\}$, the singleton containing the empty word; a nullary rule selects one symbol of $\mathcal A$.}
The other two statements are simple conclusions from Proposition~\ref{prop:uniqueness}.
\begin{prop}
\label{prop:props}
\rev{Let $n\geq0$ and let $N,N_1,N_2$ be finite subsets of $\mathbb Z$.}
Then
\begin{enumerate}[label=(\roman*)]
\item
$
L_\mathcal{A}(n) = \bigcup_{0 \le m \le n} \{ g \circ h \ | \ g \in \overline{L}_\mathcal{A}(m), h \in H_\mathcal{A}(n,m) \};
$
\item
$
G_\mathcal{A}(N) = \bigcup_{M\subseteq N} \overline{G}_\mathcal{A}(M);
$
\item
$
G_\mathcal{A}(N_1) \cap G_\mathcal{A}(N_2) = G_\mathcal{A}(N_1 \cap N_2).
$
\end{enumerate}
\end{prop}
\section{Local Rules}
\label{sec:loc}

\subsection{Symmetry Transformations}
\label{sec:loc-symm}

We now describe two transformations on the local rule space, one is based
on the permutation of the state set, the other one results from a reflection of a configuration.

Let $S_\mathcal{A}$ be the symmetric group of degree $|\mathcal{A}|$,
that is the set of all permutations of the set $\mathcal{A}$,
and suppose $\nu \in S_\mathcal{A}$.
If $a \in \mathcal{A}$ we write the image of $a$ under $\nu$ as product $\nu a$.
The extension of $\nu$ to words is defined by elementwise application.
If $w = a_1\ldots a_{n} \in \mathcal{A}^n$ is a word,
set $\nu w  = (\nu a_1) \ldots (\nu a_{n})$.
Suppose $f$ is a local rule from $\mathcal{A}^n$ to $\mathcal{A}$.
The permutation operator $\hat{\nu}$ is
defined by
$\hat{\nu}f(w) = \nu f ( \nu^{-1} w)$ for all words $w \in \mathcal{A}^n$.
Note that the  ``hat'' on the operator is necessary, because $\nu f$ and
$\hat{\nu} f$ are distinct entities.
The first one is the composite function $\nu \circ f$, whereas the second represents
the composite function $\nu \circ f \circ \nu^{-1}$.

If $w = a_1\ldots a_{n} \in \mathcal{A}^n$ is a word,
define the reflection of $w$ by $rw  = a_{n} \ldots a_1$.
Let $f$ be a local rule.
The reflection operator $\hat{r}$ is defined by $\hat{r} f(w) = f(rw)$.
Since $r$ is self-inverse ($r^{-1} = r$),
and $ra = a$ for all symbols $a \in \mathcal{A}$, we can also write $\hat{r} f(w) = rf(r^{-1}w)$,
which makes the notation consistent with that of the permutation operator.

Both operators, the permutation and the reflection operator, transform the rule
space $L_\mathcal{A}(n)$.
Because both of the operators map an irreducible local rule to another irreducible rule,
the operators can also be restricted to the rule space $\overline{L}_\mathcal{A}(n)$.
The set $R = \{1, r\}$, is called the reflection group.
The direct product of $R$ and $S_\mathcal{A}$, written as $RS_\mathcal{A}$, is the group that contains all permutations,
the reflection and their products.
We denote the group $RS_\mathcal{A}$ by  $\mathcal{S}_\mathcal{A}$.

Suppose that $\alpha$ and $\beta$ are in $\mathcal{S}_\mathcal{A}$.
Then
\[
\hat{\alpha} \hat{\beta} f (w)  = \alpha \hat{\beta} f(\alpha^{-1} w)
= \alpha \beta f( \beta^{-1} \alpha^{-1} w) = \widehat{\alpha \beta} f(w).
\]
The operators form a group that is in general isomorphic to $\mathcal{S}_\mathcal{A}$,
\rev{but for $n=0$ or $n=1$ the action is not faithful, since reflection acts trivially.}
Note that the reflection operator commutes with all permutation operators.
If $\alpha \in \mathcal{S}_\mathcal{A}$, and $f$ is a local rule such that $\hat{\alpha} f = f$,
the local rule $f$ is said to be invariant under the symmetry transformation $\alpha$.

A group action of a group $G$ on a set $A$ is a map of $G \times A$ into $A$
satisfying the following properties, see e.g.~\cite{rotman2012introduction}:
\begin{enumerate}[label=(\roman*)]
\item $g_1(g_2a) = (g_1g_2)a$ for all $g_1,g_2 \in G$, $a \in A$, and
\item $1a=a$, for all $a \in A$.
\end{enumerate}

The map $\mathcal{S}_\mathcal{A} \times L_\mathcal{A}(n) \to L_\mathcal{A}(n)$; \rev{$(\alpha, f) \mapsto \hat{\alpha} f$}
fulfils the properties of a group action.
The relation on $L_\mathcal{A}(n)$, defined by $f \sim g$ if and only if $f = \hat{\alpha} g$ for some
$\alpha \in \mathcal{S}_\mathcal{A}$, is an equivalence relation.
The equivalence classes $[f] = \{ \hat{\alpha}f \, | \, \alpha \in \mathcal{S}_\mathcal{A} \}$ form a partition
of the rule space $L_\mathcal{A}(n)$.
The set of all equivalence classes of $L_\mathcal{A}(n)$ is denoted by $L_\mathcal{A}(n) / \mathcal{S}_\mathcal{A}$.

Let $U$ be a subgroup of $\mathcal{S}_\mathcal{A}$ (denoted $U \le \mathcal{S}_\mathcal{A}$).
The conjugacy class of $U$ is the set of subgroups
$[U] = \{ V \le \mathcal{S}_\mathcal{A} \, | \, \mbox{ $V = \alpha U \alpha^{-1}$ for some $\alpha \in \mathcal{S}_\mathcal{A}$} \}$.
The stabilizer of a local rule $f \in L_\mathcal{A}(n)$ is the set
$\stab(f) = \{ \alpha \in \mathcal{S}_\mathcal{A} \ | \rev{\hat{\alpha}f = f} \}$.
An equivalence class $\rev{[f] \in L_\mathcal{A}(n) / \mathcal{S}_\mathcal{A}}$ is said to be of
type $[U]$ if the stabilizer of some $g$ in $[f]$ belongs to $[U]$.

If $\mathcal{A} = \{0,1\}$, we write $S_2$ for $S_{\{0,1\}}$, $L_2(n)$ for $L_{\{0,1\}}(n)$, and so on for the other notation
we introduced so far.
The group $S_2$ consists of the two elements $1$ and $c$, where $1$ is the identity, and $c = (01)$ is the transposition of
0 and 1.
The group $\mathcal{S}_2$ is Abelian, consists of the four elements $1$, $r$, $c$, and $rc$, and is isomorphic to the Klein four-group.
\rev{For an Abelian group, conjugacy classes of subgroups are singletons, so an orbit type can be identified with its stabilizer subgroup.}
Hence the types of $\mathcal{S}_2$ are the subgroups of $\mathcal{S}_2$ which are $\langle 1 \rangle$, $\langle r \rangle$, $\langle c \rangle$,
$\langle rc \rangle$, and $\langle r, c \rangle = \mathcal{S}_2$.
\begin{revision}
The notation $\langle\alpha,\beta,\ldots\rangle$ denotes the subgroup generated by the displayed elements.
\end{revision}

Define
\[
t^U_n = \{ [f] \in L_2(n) / \mathcal{S}_2 \, | \, \stab(f) = U \}
\]
for $U \le \mathcal{S}_2$.
The set $t^U_n$ consists of all equivalence classes of $L_2(n)$ that have type $U$.
The sets $t^U_n$, $U \le \mathcal{S}_2$, \rev{form a partition of $L_2(n) / \mathcal{S}_2$}, so
\[
L_2(n) / \mathcal{S}_2 = \bigcup_{U \le \mathcal{S}_2} t^U_n.
\]
The calculation of the number of equivalence classes
$|L_2(n) / \mathcal{S}_2|$ and of the numbers $|t^U_n|$, $U \le \mathcal{S}_2$,
can be found in \cite{SCHALLER2025100298}, where groups, group actions, and related concepts in the context of CAs
are also presented in more detail.


\subsection{Irreducible Binary Rules}

The restriction of the group action to
$\mathcal{S}_2 \times \overline{L}_2(n) \to \overline{L}_2(n)$ is well-defined.
The group action leads to the equivalence classes $\overline{L}_2(n) / \mathcal{S}_2$
and the classification by types
\[
\overline{t}^U_n = \{ [f] \in \overline{L}_2(n) / \mathcal{S}_2 \, | \, \stab(f) = U \}
\]
 for $U \le \mathcal{S}_2$.
We will now derive the number of equivalence classes of $\overline{L}_2(n) / \mathcal{S}_2$ and the numbers
$|\overline{t}^U_n|$, $U \le \mathcal{S}_2$,
from the numbers $|t^U_n|$, $U \le \mathcal{S}_2$.

For any projection $h$, we define $\hat{r} h = r \circ h \circ r$.
The projection $h$ is said to be reflection-invariant if $\hat{r}h = h$.
Let $g \circ h$ be a local rule,  where $g$ is an irreducible local rule and $h$ is a projection.
Then
\[
\hat{r}(g \circ h)  = g \circ h \circ  r  = g \circ r \circ r \circ h \circ r =
\hat{r}g \circ \hat{r}h.
\]
We have adopted the convention that a unary function operator binds more strongly than function composition.
It follows that if $f = g \circ h$, then $\hat{r} f = f$ if and only if $g = \hat{r}g$ and $h = \hat{r}h$.
Suppose $h: \mathcal{A}^n \to \mathcal{A}^m$ is a projection and $H(h) = \{p_1, \ldots, p_m\}$.
Then
\begin{align*}
\hat{r}h(a_1\ldots a_n) & = r \circ h \circ r(a_1\ldots a_n) =
r \circ h(a_n \ldots a_1) \\
& = r(a_{n+1-{p_1}} \ldots a_{n+1-{p_m}})
= a_{n+1-{p_m}} \ldots a_{n+1-p_1}.
\end{align*}
We see that the index set of $\hat{r}h$ relates to the index set of $h$ by
$H(\hat{r}h) = n + 1 - H(h)$.
We need to know how many projections $h: \mathcal{A}^n \to \mathcal{A}^m$ are reflection-invariant.
The question can be rephrased by asking how many ways there are to arrange $m$ balls symmetrically into
$n$ slots, \rev{with each slot} holding at most one ball.
\begin{prop}
\label{prop:inv-h}
Suppose there are $m$ balls and $n$ slots \rev{where $0\leq m\leq n$ and each slot can hold at most one ball.}
Then there are $\alpha_{n,m}$ arrangements of balls in slots that are reflection-symmetrical,
where $\alpha_{n,m}$ is defined by
\[
\alpha_{n,m} = \left\{
\begin{array}{ll}
0 & \mbox{if $n$ is even and $m$ is odd}, \\
\binom{\lfloor \frac{n}{2} \rfloor }{\lfloor \frac{m}{2} \rfloor } & \mbox{otherwise;}
\end{array}
\right.
\]
and where $\lfloor.\rfloor : \mathbb{R} \to \mathbb{Z}$,
$\lfloor x \rfloor = \mathrm{max}\{ p \in \mathbb{Z} \, | \, p \le x\}$, is the floor function.
The number of arrangements that are not reflection-symmetrical is denoted by $\beta_{n,m}$ defined by
\[
\beta_{n,m} = \binom{n}{m} - \alpha_{n,m}.
\]
\end{prop}
\begin{proof}
We distinguish between four cases in which $m$ is even or odd and $n$ is even or odd.

If $m$ and $n$ are even, there are $\binom{n/2}{m/2}$ ways to
arrange $m/2$ balls on the left half of the slots.
The arrangement of the balls on the right half of the slots follows from the reflection symmetry.

If $m$ is even and $n$ is odd, the slot in the center must remain empty,
so there are $\binom{(n-1)/2}{m/2}$ possibilities.

If $m$ is odd and $n$ is odd, and if $m \ge 1$, one ball must go into the slot in the middle.
Therefore, there are $\binom{(n-1)/2}{(m-1)/2}$ possibilities.

If $m$ is odd and $n$ is even, the number of balls in the left half of the slots cannot be equal
to the number of balls in the right half of the slots.
Therefore, there is no arrangement.

This proves the statement about $\alpha_{n,m}$.
Since $\alpha_{n,m} + \beta_{n,m} = \binom{n}{m}$ the statement about $\beta_{n,m}$ follows.
\end{proof}
\begin{revision}
\locmark{loc:stabilizer-proof}
Fix an irreducible local orbit $[g]\in\overline t_m^U$, and write a
padded rule as $f=g\circ h$, where $h$ is a projection from $n$ inputs
onto $m$ inputs. State complementation preserves the selected
coordinates, whereas reflection replaces $h$ by $\hat rh$. Uniqueness
of the irreducible factorization therefore gives
\begin{equation}\label{eq:projection-stabilizer}
\stab(g\circ h)=
\begin{cases}
U,&h=\hat rh,\\
U\cap\langle c\rangle,&h\ne\hat rh.
\end{cases}
\end{equation}
Indeed, an element involving reflection can fix $g\circ h$ only if it
fixes the projection's index set; once that set is fixed, it must also
fix $g$. Elements involving only state complementation fix $h$
automatically and fix $g\circ h$ exactly when they fix $g$.

Each of the $\alpha_{n,m}$ reflection-invariant projections produces
one padded orbit of type $U$. The remaining projections occur in
$\beta_{n,m}/2$ distinct pairs $\{h,\hat rh\}$. For each such pair,
the set
\[
\{g'\circ h,g'\circ\hat rh:g'\in[g]\}
\]
has $2|[g]|=8/|U|$ elements, all with stabilizer
$U\cap\langle c\rangle$. By the orbit--stabilizer theorem and
Lagrange's theorem, an orbit with stabilizer $V$ has size
$[\mathcal S_2:V]=4/|V|$. Thus each padded orbit here has
$4/|U\cap\langle c\rangle|$ elements, and the displayed set consists of
$2|U\cap\langle c\rangle|/|U|$ padded orbits. The contributions
of one irreducible orbit are thus as follows:
\[
\begin{array}{c|c|c}
U&\text{symmetric projections}&\text{asymmetric projection pairs}\\\hline
\langle1\rangle&\alpha_{n,m}\text{ of type }\langle1\rangle
 &\beta_{n,m}\text{ of type }\langle1\rangle\\
\langle c\rangle&\alpha_{n,m}\text{ of type }\langle c\rangle
 &\beta_{n,m}\text{ of type }\langle c\rangle\\
\langle r\rangle&\alpha_{n,m}\text{ of type }\langle r\rangle
 &\beta_{n,m}/2\text{ of type }\langle1\rangle\\
\langle rc\rangle&\alpha_{n,m}\text{ of type }\langle rc\rangle
 &\beta_{n,m}/2\text{ of type }\langle1\rangle\\
\mathcal S_2&\alpha_{n,m}\text{ of type }\mathcal S_2
 &\beta_{n,m}/2\text{ of type }\langle c\rangle
\end{array}
\]
Different irreducible orbits or different projection pairs give disjoint
sets of padded rules, again by unique factorization. Multiplying each
row by $|\overline t_m^U|$ and summing over $0\leq m\leq n$ gives
the unrestricted type counts. Since $\alpha_{n,n}=1$ and
$\beta_{n,n}=0$, isolating the term $m=n$ gives the following recurrences.
\end{revision}

\begin{prop}
\rev{\label{prop:recurrences}\locmark{loc:recurrences}The numbers of equivalence classes of $\overline L_2(n)$ by type are given by}
\begin{equation*}
  \relax
  \begin{aligned}
|\overline{t}_n^{\mathcal{S}_2}| &= |t_n^{\mathcal{S}_2}| - \sum_{m=0}^{n-1} \alpha_{n,m} |\overline{t}_m^{\mathcal{S}_2}|; \\
|\overline{t}_n^{\langle r \rangle}| &= |t_n^{\langle r \rangle}|
        - \sum_{m=0}^{n-1} \alpha_{n,m} |\overline{t}_m^{\langle r \rangle}|; \\
|\overline{t}_n^{\langle rc \rangle}| &= |t_n^{\langle rc \rangle}|
        - \sum_{m=0}^{n-1} \alpha_{n,m} |\overline{t}_m^{\langle rc \rangle}|; \\
|\overline{t}_n^{\langle c \rangle}| &= |t_n^{\langle c \rangle}|
        - \sum_{m=0}^{n-1} \left[ \binom{n}{m} |\overline{t}_m^{\langle c \rangle}|
        + \frac{1}{2} \beta_{n,m} |\overline{t}_m^{\mathcal{S}_2}| \right]; \\
|\overline{t}_n^{\langle 1 \rangle}| &=  |t_n^{\langle 1 \rangle}|
        - \sum_{m=0}^{n-1} \left[ \binom{n}{m} |\overline{t}_m^{\langle 1 \rangle}|
        + \frac{1}{2} \beta_{n,m}  \left( |\overline{t}_m^{\langle r \rangle}|
        + |\overline{t}_m^{\langle rc \rangle}| \right) \right].
  \end{aligned}
\end{equation*}
\end{prop}

\begin{table}
\begin{center}
\caption{Number of equivalence classes of irreducible local rules.}
\label{tab:s2_fs}
\begin{tabular}{|l|r|r|r|r|r|r|}
\hline
       & \multicolumn{6}{|c|}{\rev{$|\overline{t}^U_n|$}} \\ \hline
\rev{$U$}    & $n=0$ & $n=1$ & $n=2$ & $n=3$ & $n=4$ & $n=5$ \\ \hline
\rev{$\mathcal S_2$}                   & 0 & 2 & 0 &  6 &       0 & 1\,010 \\ \hline
$\langle rc \rangle$                            & 0 & 0 & 0 &  4 &       0 & 32\,248 \\ \hline
$\langle c \rangle$                             & 0 & 0 & 0 &  2 &     104 & \rev{31\,688} \\ \hline
$\langle r \rangle$                             & 1 & 0 & 3 & 24 &     505 & 523\,216 \\ \hline
$\langle 1 \rangle$                             & 0 & 0 & 1 & 38 & 15\,844 & 1\,073\,336\,680 \\ \hline
$|\overline{L}_2(n)/\mathcal{S}_2|$     & 1 & 2 & 4 & 74 & 16\,453 & 1\,073\,954\,842 \\ \hline
\end{tabular}
\end{center}
\end{table}

The numbers \rev{$|\overline{t}^U_n|$} \rev{can now be calculated iteratively, starting at $n=0$.}
Table~\ref{tab:s2_fs} depicts the results for a neighbourhood size $n=0, \ldots, 5$.
The sets $\overline{t}^U_n$ form a partition of $\overline{L}_2(n) / \mathcal{S}_2$, so
\[
|\overline{L}_2(n) / \mathcal{S}_2| = \sum_{U \le \mathcal{S}_2} |\overline{t}^U_n|.
\]


\subsection{Small Neighbourhoods}

We still limit ourselves to binary CAs.
There are only seven equivalence classes for a neighbourhood size less than or equal to two, which we
list below.

\begin{enumerate}

\item Cardinality Zero, $\overline{L}_2(0)$.
These are the constant functions $W^0_0 = 0$ and $W^0_1 = \hat{c} W^0_0 = 1$.

\item Cardinality One, $\overline{L}_2(1)$.
One function is the identity $W^1_2 = \id$, $W^1_2(a_0) = a_0$ for $a_0 \in \{0,1\}$.
The second function inverts the $\id$ function, $W^1_1(a_0) = ca_0$, $W^1_1 = c \id$.
Note that $\hat{c} \id = c \circ \id \circ c = \id \ne c \circ \id$.

\item Cardinality Two, $\overline{L}_2(2)$.

Equivalence Class $\{W^2_1, W^2_7\}$.
The two functions are $W^2_1$, given by
\[
W^2_1(a_0a_1) =
\left\{
\begin{array}{ll}
1 & \mbox{if $a_0a_1 = 00$;} \\
0 & \mbox{otherwise,}
\end{array}
\right.
\]
and $W^2_{7} = \hat{c}W^2_1$.

Equivalence Class $\{W^2_6, W^2_9\}$.
The two functions are
$W^2_6$, given by
\[
W^2_6(a_0a_1) =
\left\{
\begin{array}{ll}
1 & \mbox{if $a_0a_1 = 01$ or $a_0a_1 = 10$;} \\
0 & \mbox{otherwise,}
\end{array}
\right.
\]
and $W^2_9 = \hat{c}W^2_6$.

Equivalence Class $\{W^2_8, W^2_{14}\}$.
The two functions are
$W^2_8$, given by
\[
W^2_8(a_0a_1) =
\left\{
\begin{array}{ll}
1 & \mbox{if $a_0a_1 = 11$;} \\
0 & \mbox{otherwise,}
\end{array}
\right.
\]
and $W^2_{14} = \hat{c}W^2_8$.

Equivalence Class
$\{W^2_2,W^2_4,W^2_{11},W^2_{13}\}$.
The first function is given by
\[
W^2_2(a_0a_1) =
\left\{
\begin{array}{ll}
1 & \mbox{if $a_0a_1 = 01$;} \\
0 & \mbox{otherwise\rev{.}}
\end{array}
\right.
\]
The other ones are $W^2_4 = \hat{r}W^2_2$, $W^2_{11} = \hat{c}W^2_2$, and
$W^2_{13} = \widehat{rc} W^2_2$.


\end{enumerate}

\section{Global Maps}
\label{sec:glob}

\subsection{Symmetry Transformations}
We now examine the symmetry transformations on global maps and link two of them to the
symmetry transformations on local rules, see  Subsection \ref{sec:loc-symm}.

First we consider the permutation operator.
Let $S_\mathcal{A}$ be the symmetric group of degree $|\mathcal{A}|$
and suppose $\nu \in S_\mathcal{A}$.
The extension of $\nu$ to configurations is defined by elementwise application.
If $x$ is a configuration, set $(\nu x)_i = \nu (x_i)$.
If $\Phi^N_f$ is the induced global mapping of $f$, we define
$\hat{\nu} \Phi^N_f$ by
$\hat{\nu} \Phi^N_f (x) = \nu \Phi^N_f (\nu^{-1} x)$ for all configurations $x$.
From
\begin{eqnarray*}
\left( \hat{\nu} \Phi^N_f(x) \right)_{i}
& = & \left( \nu \Phi^N_f (\nu^{-1} x) \right)_{i}
= \nu f \left(\nu^{-1}(x_{i+j_1} \ldots x_{i+j_{n}}) \right) \\
& = & \hat{\nu} f(x_{i+j_1} \ldots x_{i+j_{n}})
 =  \Phi^N_{\hat{\nu} f} (x)_{i}
\end{eqnarray*}
follows $\hat{\nu} \Phi^N_f =  \Phi^N_{\hat{\nu} f}$.
Let $A = (\mathcal{A}, N, f)$ be a CA and put $B = (\mathcal{A}, N, \hat{\nu} f)$.
By defining $\hat{\nu} A = B$, the permutation operator $\hat{\nu}$ maps the CA $A$ to the CA $B$.

The two other transformations, reflection and shift, originate from isometric transformations on the set of integers.
Isometric transformations, that is transformations that preserve the distance, form a group,
the infinite dihedral group $D_\infty$, see e.g. \cite{armstrong1997groups}.
By definition if $\varphi \in D_\infty$ and $i,j \in \mathbb{Z}$,
then $|\varphi(i) - \varphi(j)| = |i - j|$.
The group $D_\infty$ has the representation $\langle r, t \, | \, r^2=1, rtr = t^{-1} \rangle$,
where $r$ and $t$ can be chosen to be $r(i) = -i$ and $t(i) = i+1$ for all $i \in \mathbb{Z}$.

We extend these operations to sets of integers.
If $N$ is any set of integers, put $-N = \{-i \, | \, i \in N \}$ and
$j + N = \{ j + i \, | \, i \in N \}$ if $j$ is itself an integer.
Two sets of integers $M$ and $N$ are called congruent, in symbols $M \cong N$,  if and only if
there is an integer $j$ such that $N = j + M $ or $N = j - M$.
The congruence relation is an equivalence relation.
A set $N$ of integers is called reflection-symmetrical, if
there is an integer $j$ such that $N = j - N$.
Let $\Sym$ denote the set of all reflection-symmetrical subsets of integers.
Then $N$ is reflection-symmetrical if and only if $N \in \Sym$.
If $N \in \Sym$, then $[N] \subset \Sym$.

The shift operator $\sigma$ plays a dual role.
First, by the \rev{Curtis--Hedlund--Lyndon} Theorem the shift operator characterizes the continuous functions on $\mathcal{A}^\mathbb{Z}$
that can be expressed as the global maps of a CA.
Second it can be used to map a CA to another one.
It is this second meaning that we now study.
Let $A=(\mathcal{A},N,f)$ be a CA.
Then
\[
(\sigma \Phi^N_f(x))_i = (\Phi^N_f(x))_{i+1}
= f(x_{i+j_1+1} \ldots x_{i+j_{n}+1}) = ( \Phi^{1+N}_f(x))_i.
\]
Hence we also can see $\sigma$ as the operator that maps the CA $A = (\mathcal{A},N,f)$ to
the CA $B = \sigma A = (\mathcal{A}, 1+N,f)$,
expressed by
\[
\sigma \Phi^A = \Phi^{\sigma A} \ \ \ \mbox{or} \ \ \ \sigma \Phi^N_f = \Phi^{1+N}_f.
\]

The reflection operator is based on the reflections of configurations.
If $x$ is a configuration, set $(rx)_i = x_{-i}$.
The reflection operator $\hat{r}$ is formally defined by  \rev{$\hat{r}\Phi(x) = r\Phi(rx)$}.
Let $A = (\mathcal{A},N,f)$ be a CA with the global mapping $\Phi^N_f$.
Then
\begin{align*}
(\hat{r} \Phi^N_f(x))_i & =  (r\Phi^N_f(rx))_i
        =  (\Phi^N_f(rx))_{-i} = f((rx)_{-i+j_1} \ldots (rx)_{-i+j_n}) \\
& =  f(x_{i-j_1} \ldots x_{i-j_n}) = f(r(x_{i-j_n} \ldots x_{i-j_1})) \\
& =  \hat{r} f(x_{i-j_n} \ldots x_{i-j_1}) = (\Phi^{-N}_{\hat{r}f}(x))_i.
\end{align*}
If we define $\hat{r}A = (\mathcal{A}, -N, \hat{r}f)$, we can express the relation above as
$\hat{r} \Phi^A = \Phi^{\hat{r} A}$.
Let $x$ be a configuration.
From
\[
(\sigma r (x))_i = (r(x))_{i+1} = x_{-i-1} = (\sigma^{-1}(x))_{-i} = (r \sigma^{-1}(x))_i
\]
follows the commutator relation $\rev{\sigma r = r\sigma^{-1}}$.
Then
\[
\sigma \hat{r} \Phi^N_f = \sigma r \Phi^N_f r = r \sigma^{-1} \Phi^N_f r = \hat{r} \sigma^{-1} \Phi^N_f
\]
also implies $\sigma \hat{r} = \hat{r} \sigma^{-1}$.


In general, the reflection operator is not a transformation on the global rule space $G_\mathcal{A}(N)$ or
$\overline{G}_\mathcal{A}(N)$,
but a map $\hat{r}: G_\mathcal{A}(N) \to G_\mathcal{A}(-N)$ or $\hat{r}: \overline{G}_\mathcal{A}(N) \to
\overline{G}_\mathcal{A}(-N)$.
If $N = -N$, the map becomes a transformation.

\begin{revision}
\begin{prop}\label{prop:intrinsic-action}\locmark{loc:action}
State relabeling, reflection, and translation of local representations
induce well-defined bijections of $\mathcal C_{\mathcal A}$.
Explicitly, if $\Phi_f^N=\Phi_g^M$, then for every
$\nu\in S_{\mathcal A}$ and $a\in\mathbb Z$,
\[
\Phi_{\hat\nu f}^{N}=\Phi_{\hat\nu g}^{M},\qquad
\Phi_{\hat r f}^{-N}=\Phi_{\hat r g}^{-M},\qquad
\Phi_f^{a+N}=\Phi_g^{a+M}.
\]
These induced transformations are respectively
$\hat\nu\Phi=\nu\Phi\nu^{-1}$, $\hat r\Phi=r\Phi r$, and
$\Phi\mapsto\sigma^a\Phi$.
For the essential pair $(E,q)$ of $\Phi$, their essential pairs are
\[
(E,\hat\nu q),\qquad (-E,\hat r q),\qquad (a+E,q).
\]
\end{prop}
\begin{proof}
The coordinate identities established above hold for every local
representation, whether or not its rule is irreducible:
\[
\Phi_{\hat\nu f}^{N}=\nu\Phi_f^N\nu^{-1},\qquad
\Phi_{\hat r f}^{-N}=r\Phi_f^N r,\qquad
\Phi_f^{a+N}=\sigma^a\Phi_f^N.
\]
Replacing $\Phi_f^N$ on the right by any equal representation
$\Phi_g^M$ proves the three claimed equalities. Thus the transformed
global maps are independent of the chosen neighbourhood and local rule,
even when the two representations have different arities.
Applying the same identities to the essential pair gives the displayed
transformed pairs. Relabeling and reversal preserve irreducibility,
and translation leaves the local rule unchanged; hence these pairs
are essential by Proposition~\ref{prop:uniqueness}. The argument also
applies to nullary rules and the empty neighbourhood.
The inverse operators are $\widehat{\nu^{-1}}$, $\hat r$, and
postcomposition by $\sigma^{-a}$.
\end{proof}
In expressions such as $\sigma^a\hat r\Phi$, the shift acts by
postcomposition on the global map. On rule space the identities are
\[
\hat r^2=1,\qquad
\sigma^a\hat r=\hat r\sigma^{-a},\qquad
\hat\nu\sigma^a=\sigma^a\hat\nu,
\qquad \hat\nu\hat r=\hat r\hat\nu.
\]
These rule-space operations are distinct from the configuration maps used
to define them.
\end{revision}

The products of all symmetry transformations form the infinite group
\[
\mathcal{T}_\mathcal{A} = \{ \sigma^i \hat{r}^j \hat{\nu} \, | \, i \in \mathbb{Z}; j=0,1; \nu \in S_\mathcal{A} \}.
\]
\rev{It acts as a group on} $\mathcal{C}_\mathcal{A}$, the set of all global maps induced by CAs over the
alphabet $\mathcal{A}$.
The group action partitions $\mathcal{C}_\mathcal{A}$ into the
equivalence classes $\mathcal{C}_\mathcal{A} / \mathcal{T}_\mathcal{A}$.
Suppose the global maps $\Phi^A$ and $\Phi^B$ are in $\mathcal{C}_\mathcal{A}$.
We write $\Phi^A \cong \Phi^B$ if and only if there is an equivalence class
$C \in  \mathcal{C}_\mathcal{A} / \mathcal{T}_\mathcal{A}$ such that
both $\Phi^A$ and $\Phi^B$ are members of $C$.
Being an equivalence relation, the restriction of the relation $\cong$ to
any subset $\mathcal{D}$ of $\mathcal{C}_\mathcal{A}$ is still an equivalence relation.
We then write $\mathcal{D} / {\cong}$ for the resulting equivalence classes.
Before we study the general case, we examine the equivalence classes of global functions induced by
a neighbourhood of size $n=0,1,2$.

\begin{enumerate}

\item Let $n=0$. Then $N=\emptyset$.
\rev{The nullary rules have domain $\mathcal A^0=\{\varepsilon\}$ and are constant.}
If $f$ is a constant local map, then the equivalence class of \rev{$\Phi_f^\emptyset$} is \rev{the set
of all constant global maps on $\mathcal{A}^\mathbb{Z}$.}
Thus $[\Phi^\emptyset_f] = \{ \Phi^\emptyset_g \, | \, g() \in \mathcal{A} \}$.
There is one equivalence class and the size of this class is $k$.

\item Let $n=1$, $\rev{N=\{0\}}$, and $f \in \overline{L}_\mathcal{A}(1)$.
Then $[\rev{\Phi_f^{\{0\}}}] = \{ \rev{\Phi_g^{\{i\}}} \, | \, i \in \mathbb{Z}, g \in [f] \}$.
There are \rev{$|\overline L_\mathcal{A}(1) / \mathcal{S}_{\mathcal A}|$} equivalence classes of the form $[\rev{\Phi_f^{\{0\}}}]$,
and each equivalence class is countably infinite.

\item \rev{\locmark{loc:two-site}Let $n=2$, $p>0$, $N=\{0,p\}$,} and $f \in \overline{L}_\mathcal{A}(2)$.
From
$\rev{\sigma^p}\hat{r}\Phi^{\{0,p\}}_f = \rev{\sigma^p}\Phi^{\{-p,0\}}_{\hat{r}f}
= \Phi^{\{0,p\}}_{\hat{r}f}$
follows $\Phi^{\{0,p\}}_{\hat{r}f} \in [\Phi^{\{0,p\}}_f]$.
We conclude that $[\Phi^{\{0,p\}}_f] = \{ \Phi^{\{i,i+p\}}_g \, | \, i \in \mathbb{Z}, g \in [f] \}$.
\rev{For each fixed $p>0$, there are $|\overline L_{\mathcal A}(2)/\mathcal S_{\mathcal A}|$ such classes, each countably infinite. Allowing $p$ to vary yields countably infinitely many classes.}

\end{enumerate}

We consider now the general case.
Let $N$ be a finite set of integers and $f \in \overline{L}_\mathcal{A}(|N|)$.

\begin{enumerate}

\item Suppose $N \in \Sym$.
Then $N = p -N$ \rev{for an integer $p$}, and consequently $\Phi^{N}_{\hat{r}f} \in [\Phi^{N}_f]$.
Thus $[\Phi^N_f] = \{\Phi^{d+N}_g \, | \, d \in \mathbb{Z}, g \in [f] \}$.
Since $N \in \Sym$, the equivalence class can also be written as
$[\Phi^N_f] = \{\Phi^{M}_g \, | \, M \cong N, g \in [f] \}$.

\item Suppose $N \not\in \Sym$ and suppose there is a $\nu \in S_\mathcal{A}$ such
that $\widehat{r\nu}f = f$.
Then $\widehat{r\nu} \Phi^N_f = \Phi^{-N}_f$.
This implies as above that
$[\Phi^N_f] = \{\Phi^{M}_g \, | \, M \cong N, g \in [f] \}$.

\item Suppose $N \not\in \Sym$ and suppose there is no $\nu \in S_\mathcal{A}$
such that $\widehat{r\nu}f =f$.
Then $[\Phi^N_f] = \{ \Phi^{d+N}_{\hat{\nu}f},  \Phi^{d-N}_{\widehat{r\nu}f} \, | \, d \in \mathbb{Z},
\nu \in S_\mathcal{A} \}$.
Note that $\Phi^N_{\hat{r}f} \not\in [\Phi^N_f]$.

\end{enumerate}
By grouping the last two cases together, we obtain the following proposition.
\begin{prop}
\label{prop:global-eq}
Let $N$ be a finite set of integers and $f \in \overline{L}_\mathcal{A}(|N|)$.
\begin{enumerate}[label=(\roman*)]
\item
\label{prop:global-eq-1}
If $N \in \Sym$ or $\widehat{r\nu}f = f$ for some $\nu \in S_\mathcal{A}$ then
$N$ and $[f]$ induce the following equivalence class of $\mathcal{C}_\mathcal{A} / \mathcal{T}_\mathcal{A}$:
\[
[\Phi^N_f] = \{\Phi^{M}_g \, | \, M \cong N, g \in [f] \}.
\]
\item
\label{prop:global-eq-2}
Otherwise $N$ and $[f]$ induce the following two equivalence classes of $\mathcal{C}_\mathcal{A} / \mathcal{T}_\mathcal{A}$:
\begin{align*}
[\Phi^N_f] & = \{ \Phi^{d+N}_{\hat{\nu}f},  \Phi^{d-N}_{\widehat{r\nu}f} \, | \, d \in \mathbb{Z},
\nu \in S_\mathcal{A} \}; \\
[\Phi^N_{\hat{r}f}] & = \{ \Phi^{d+N}_{\widehat{r\nu}f},  \Phi^{d-N}_{\hat{\nu}f} \, | \, d \in \mathbb{Z},
\nu \in S_\mathcal{A} \}.
\end{align*}
\end{enumerate}
\end{prop}

\subsection{Finite Intervals}

This subsection deals with the equivalence classes of global maps or CAs whose neighbourhood
is a subset of a given finite integer interval.
We take a step back and first consider the equivalence relation based only on
permutations of the alphabet and reflection of the configuration.
Let $p$ be a nonnegative integer and consider the interval $N=[-p,p]$.
The map
\[
\mathcal{S}_\mathcal{A}  \times L_\mathcal{A}(2p+1) \to L_\mathcal{A}(2p+1); \
(\alpha, f) \mapsto \hat{\alpha}f
\]
is a group action, and so
is the map
\[
\mathcal{S}_\mathcal{A} \times G_\mathcal{A}(N) \to G_\mathcal{A}(N); \
(\alpha, \Phi^N_f) \mapsto \hat{\alpha} \Phi^N_f.
\]
We see that the map
\[
L_\mathcal{A}(2p+1) / \mathcal{S}_\mathcal{A} \to G_\mathcal{A}(N)/\mathcal{S}_\mathcal{A}; \
[f] \mapsto [\Phi^N_f]
\]
 is a bijection.
The same group acts both on the set of local maps and the set of global maps,
and the equivalence classes associated with the group actions are in one-to-one correspondence.

The size of $[-p,p]$ is always odd.
To deal with all neighbourhood sizes we now investigate the intervals $N = [1,n]$.
If $f \in L_\mathcal{A}(n)$ and $f$ is not constant, then
\[
\hat{r}\Phi^{N}_f = \Phi^{-N}_{\hat{r}f} \not \in G_\mathcal{A}(N).
\]
This means that we cannot use the same group action as above.
We note however that
\[
\sigma^{n+1}\hat{r}\Phi^N_f = \Phi^{N}_{\hat{r}f} \in G_\mathcal{A}(N).
\]
Let us therefore define the group
\[
\hat{\mathcal{S}}_\mathcal{A}(n+1) = \rev{\langle \sigma^{n+1}\hat{r}, \hat{\nu} \, | \, \nu \in S_\mathcal{A} \rangle}.
\]
From $(\sigma^{n+1}\hat{r})^{-1} = \sigma^{n+1}\hat{r}$, we conclude that
$\hat{\mathcal{S}}_\mathcal{A}(n+1)$ is isomorphic to $\mathcal{S}_\mathcal{A}$.
It follows, that
\[
\hat{\mathcal{S}}_\mathcal{A}(n+1) \times G_\mathcal{A}(N) \to G_\mathcal{A}(N); \
(\hat{\alpha}, \Phi^N_f) \mapsto \hat{\alpha} \Phi^N_f
\]
is a group action and we conclude that
$L_\mathcal{A}(n) / \mathcal{S}_\mathcal{A}$
and $G_\mathcal{A}(N)/\hat{\mathcal{S}}_\mathcal{A}(n+1)$ can be again put in one-to-one correspondence.

Let us now include the shift operator and let us consider the group $\mathcal{T}_\mathcal{A}$
of symmetry operators.
Without loss of generality we now consider the intervals $[1,n]$;
the counting is identical for the intervals of the form $[p, p+n-1]$.
To start with, we present the following example.
We ask which equivalence classes of $G_2([1,4])/{\cong}$ are induced
by some of the local rules of $\overline{L}_2(3)$ and neighbourhoods that are subsets of the interval $[1,4]$.
There are  four 3-element subsets of the integer interval $[1,4]$, which we denote by
$A=[1,3]$, $B=[2,4]$, $C=\{1,2,4\}$, and $D=\{1,3,4\}$.
Let us now consider the equivalence classes
$[W_{54}]$, $[W_{57}]$, $[W_{110}] \in L_2(3) / \mathcal{S}_2$.
The types of these equivalence classes are $\langle r \rangle$, $\langle rc \rangle$, and $\langle 1 \rangle$, respectively.
Note that $\sigma^4 \hat{r} \Phi^A_f = \Phi^A_{\hat{r}f}$, $\sigma \Phi^A_f = \Phi^B_f$, and
$\sigma^5 \hat{r} \Phi^C_f = \Phi^D_{\hat{r}f}$.
\begin{revision}
In the following display, brackets denote classes restricted to $G_2([1,4])$;
the ambient classes also contain translates outside this interval.
\end{revision}
\begin{align*}
[\Phi^{A}_{W_{54}}] & = \{ \Phi^{A}_{W_{54}}, \Phi^{A}_{\hat{c}W_{54}}, \Phi^{B}_{W_{54}},
        \Phi^{B}_{\hat{c}W_{54}} \}; \\
[\Phi^{A}_{W_{57}}] & = \{ \Phi^{A}_{W_{57}}, \rev{\Phi^{A}_{\hat{r}W_{57}}}, \Phi^{B}_{W_{57}},
        \Phi^{B}_{\hat{r}W_{57}} \}; \\
[\Phi^{A}_{W_{110}}] & = \{ \Phi^{A}_{W_{110}}, \Phi^{A}_{\hat{c}W_{110}}, \Phi^{A}_{\hat{r}W_{110}},
        \Phi^{A}_{\widehat{rc}W_{110}},
        \Phi^{B}_{W_{110}}, \Phi^{B}_{\hat{c}W_{110}}, \Phi^{B}_{\hat{r}W_{110}},
        \Phi^{B}_{\widehat{rc}W_{110}}
\}; \\
[\Phi^{C}_{W_{54}}] & = \{ \Phi^{C}_{W_{54}}, \Phi^{C}_{\hat{c}W_{54}}, \Phi^{D}_{W_{54}},
        \Phi^{D}_{\hat{c}W_{54}} \}; \\
[\Phi^{C}_{W_{57}}] & = \{ \Phi^{C}_{W_{57}}, \Phi^{C}_{\hat{r}W_{57}}, \Phi^{D}_{W_{57}},
        \Phi^{D}_{\hat{r}W_{57}} \}; \\
[\Phi^{C}_{W_{110}}] & = \{ \Phi^{C}_{W_{110}}, \Phi^{C}_{\hat{c}W_{110}}, \Phi^{D}_{\hat{r}W_{110}},
        \Phi^{D}_{\widehat{rc}W_{110}} \}; \\
[\Phi^{C}_{\hat{r}W_{110}}] & = \{ \Phi^{C}_{\hat{r}W_{110}}, \Phi^{C}_{\widehat{rc}W_{110}},
        \Phi^{D}_{W_{110}}, \Phi^{D}_{\hat{c}W_{110}} \}.
\end{align*}
The neighbourhood $A$ and each of the equivalence classes $[W_{110}]$, $[W_{54}]$, and $[W_{57}]$
induce exactly one equivalence class of global maps in $G_2([1,4])/{\cong}$.
The situation is different for the neighbourhood $C$.
In this case, $C$ and each of $[W_{54}]$ and $[W_{57}]$ induce one equivalence class of global maps,
but $C$ and $[W_{110}]$ induce two.
We explain this by observing that $\hat{r}W_{110} \ne W_{110}$ and
$\hat{r}W_{110} \ne \hat{c}W_{110}$.
This example demonstrates that the equivalence classes of $G_2([1,4])/{\cong}$ can be determined by examining
the equivalence classes of $\rev{\overline{L}_2(n)/\mathcal{S}_2}$, $0 \le n \le 4$, and the neighbourhoods contained
in $[1,4]$, while taking into account the types of these equivalence classes.

Earlier on, we defined the relation $\cong$ on the subsets of integers.
We restrict this relation now to the power set of a given subset $N$ of integers, and denote the equivalence
classes by $\mathcal{P}(N)/{\cong}$.
Thus $\mathcal{P}(N)/{\cong}$ is the set $ \{ [M] \, | \, M \subseteq N \}$, where $[M]$ is the equivalence
class $\{ M^\prime \subseteq N \, | \, M^\prime \cong M \}$.
We are using this notation in the following proposition.
\begin{revision}
\begin{prop}\label{prop:global-shift}\locmark{loc:congruence-count}
Let $N=[1,n]$. The number of congruence classes is
\begin{equation}\label{eq:congruence-count}
\begin{split}
|G_2(N)/{\cong}|={}&
\sum_{[M]\in\mathcal P(N)/{\cong}}
  |\overline L_2(|M|)/\mathcal S_2|\\
&+\sum_{\substack{[M]\in\mathcal P(N)/{\cong}\\ M\notin\Sym}}
 \left(|\overline t_{|M|}^{\langle1\rangle}|
       +|\overline t_{|M|}^{\langle c\rangle}|\right).
\end{split}
\end{equation}
\end{prop}
\begin{proof}
Minimal-memory uniqueness gives the disjoint decomposition
\[
G_2(N)=\bigsqcup_{M\subseteq N}
       \{\Phi_f^M:f\in\overline L_2(|M|)\}.
\]
Fix a congruence class of supports, with representative $M$ and $m=|M|$.
If $M\in\Sym$, Proposition~\ref{prop:global-eq}(i) associates exactly
one global congruence class to each local orbit in
$\overline L_2(m)/\mathcal S_2$.
If $M\notin\Sym$, the same is true for the local types
$\langle r\rangle$, $\langle rc\rangle$, and $\mathcal S_2$.
For the two remaining types, $\langle1\rangle$ and
$\langle c\rangle$, Proposition~\ref{prop:global-eq}(ii) gives two
distinct global classes. Both meet $G_2(N)$, because both $f$ and
$\hat r f$ can be placed on $M$ itself. Thus an asymmetric support
class contributes
\[
|\overline L_2(m)/\mathcal S_2|
+|\overline t_m^{\langle1\rangle}|
+|\overline t_m^{\langle c\rangle}|.
\]
Summing over the support classes proves the formula. Every stabilizer
count here is a count of irreducible rules.
\end{proof}
\end{revision}

The equivalence classes of $\mathcal{P}([1,n])/{\cong}$ for small values of $n$ are given by
\begin{align*}
\mathcal{P}([1,0]) & : [\emptyset] \\
\mathcal{P}([1,1]) & : [\emptyset], [\{1\}];  \\
\mathcal{P}([1,2]) & : [\emptyset], [\{1\}],  [\{1,2\}]; \\
\mathcal{P}([1,3]) & : [\emptyset], [\{1\}],  [\{1,2\}], [\{1,3\}], [\{1,2,3\}]; \\
\mathcal{P}([1,4]) & : [\emptyset], [\{1\}],  [\{1,2\}], [\{1,3\}], [\{1,4\}], [\{1,2,3\}], [\{1,2,4\}], [\{1,2,3,4\}]; \\
\mathcal{P}([1,5]) & : [\emptyset], [\{1\}],  [\{1,2\}], [\{1,3\}], [\{1,4\}], [\{1,5\}], [\{1,2,3\}], [\{1,2,4\}], [\{1,2,5\}], \\
        &  [\{1,3,5\}], [\{1,2,3,4\}], [\{1,2,3,5\}], [\{1,2,4,5\}], [\{1,2,3,4,5\}].
\end{align*}
From these equivalence classes, the following are not reflection-symmetrical: $[\{1,2,4\}]$, $[\{1,2,5\}]$, and $[\{1,2,3,5\}]$.
Therefore by Proposition~\ref{prop:global-shift} and Table~\ref{tab:s2_fs}
\begin{align}
|G_2([1,0])/{\cong}| & = |\overline{L}_2(0) / \mathcal{S}_2| = 1; \notag \\
|G_2([1,1])/{\cong}| & = |\overline{L}_2(0) / \mathcal{S}_2| + |\overline{L}_2(1) / \mathcal{S}_2| = 3; \notag \\
|G_2([1,2])/{\cong}| & = |\overline{L}_2(0) / \mathcal{S}_2| + |\overline{L}_2(1) / \mathcal{S}_2|
+ |\overline{L}_2(2) / \mathcal{S}_2| = 7; \notag \\
|G_2([1,3])/{\cong}| & = |\overline{L}_2(0) / \mathcal{S}_2| +  |\overline{L}_2(1) / \mathcal{S}_2|
+ 2 |\overline{L}_2(2) / \mathcal{S}_2|
 + |\overline{L}_2(3) / \mathcal{S}_2| = 85; \label{eq:classes85}\\
|G_2([1,4])/{\cong}| & = |\overline{L}_2(0) / \mathcal{S}_2| + |\overline{L}_2(1) / \mathcal{S}_2|
+ 3 |\overline{L}_2(2) / \mathcal{S}_2| + 2 |\overline{L}_2(3) / \mathcal{S}_2| \notag \\
& + |\overline{t}^{\langle 1 \rangle}_3| + |\overline{t}^{\langle c \rangle}_3|
+ |\overline{L}_2(4) / \mathcal{S}_2| = 16\,656; \notag \\
|G_2([1,5])/{\cong}| & = |\overline{L}_2(0) / \mathcal{S}_2| + |\overline{L}_2(1) / \mathcal{S}_2|
+ 4 |\overline{L}_2(2) / \mathcal{S}_2|+ 4 |\overline{L}_2(3) / \mathcal{S}_2| \notag \\
& + 2|\overline{t}^{\langle 1 \rangle}_3| + 2|\overline{t}^{\langle c \rangle}_3|
+ 3 |\overline{L}_2(4) / \mathcal{S}_2| + |\overline{t}^{\langle 1 \rangle}_4| + |\overline{t}^{\langle c \rangle}_4| \notag \\
& +  |\overline{L}_2(5) / \mathcal{S}_2| = 1\,074\,020\,544 \notag.
\end{align}
\begin{revision}
\paragraph{Independent verification of the four-site count.}
\locmark{loc:four-site}
For $n=4$, the support classes represented after translation by
\[
\emptyset,\quad\{0\},\quad\{0,1\},\quad\{0,2\},\quad\{0,3\},
\quad\{0,1,2\},\quad\{0,1,3\},\quad\{0,1,2,3\}
\]
contribute respectively
\[
1,\quad2,\quad4,\quad4,\quad4,\quad74,\quad114,\quad16\,453.
\]
The asymmetric three-site contribution is
$114=74+38+2$, giving
\begin{equation}\label{eq:four-site-audit}
|G_2([1,4])/{\cong}|
=1+2+3\cdot4+74+(74+38+2)+16\,453
=16\,656.
\end{equation}
This value is independently obtained by enumerating all $65\,536$
four-input Boolean truth tables and canonicalizing their essential pairs;
Appendix~\ref{app:verification} specifies the algorithm.
In particular, reflection must transform the support and the reduced rule
together. Reflecting the rule alone on an asymmetric support would
incorrectly merge distinct classes.

\paragraph{Comparison with Ruivo et al.}
\locmark{loc:ref16}
Table~1 on p.~287 of \cite{RUIVO2018280} reports $16\,032$ extended
dynamical equivalence classes (DECx) for the binary radius-$3/2$ space
$\mathcal R_{(2,3/2)}$. Its four-site neighbourhood is $[-1,2]$, whose
translation to $[1,4]$ does not change the count. Definitions~6 and~7 of
that paper combine shifts, reversal, and state complementation and give
the same relation as $\cong$ restricted to a fixed interval. In particular,
any nonconstant rule whose translated essential support still lies in a
four-site interval has shift displacement at most three, exactly their
bound $2r=3$. Allowing ambient translates here therefore makes no extra
identifications within that interval.

The difference is a numerical error in that DECx table entry, not a
difference between the two equivalence relations. Two checks make this
conclusion independent of assumptions about the earlier implementation.
First, fully essential four-input rules alone form $16\,453$ local
symmetry classes. Such a support admits no nonzero translation within
the four-site interval, so these classes remain distinct under DECx.
Consequently $16\,032$ is already below the necessary lower bound.
Second, the same published row gives $65\,282$ shift classes (SEC)
among $65\,536$ rules and $16\,704$ ordinary dynamical classes (DEC).
Joining the shift partition to any partition can reduce the number of
classes by at most $65\,536-65\,282=254$: connect each shift class by
a spanning tree and add its edges to the existing partition.
Thus the published SEC and DEC counts imply
\[
\#\mathrm{DECx}\geq16\,704-254=16\,450>16\,032.
\]

A second exhaustive program implements the shifts of Definitions~6--7
directly on the padded four-input truth tables, without reducing supports
or using the counting formula. It reproduces the published SEC count
$65\,282$ and DEC count $16\,704$, and obtains DECx $16\,656$.
There are 46 mergers of two DEC classes and one merger of three, hence
$16\,704-46-2=16\,656$.
The program and the list of all these mergers are supplied with the
verification package. These independent computations and the support
decomposition \eqref{eq:four-site-audit} establish the corrected value;
they do not determine whether the erroneous printed entry arose in
calculation or transcription.
\end{revision}


\begin{revision}
\subsection{Scaling}
\label{sec:scaling}
Scaling changes distances between essential sites while retaining the
ordered local rule. Its definition must depend on the global map itself,
rather than on a neighbourhood padded with unused sites.

\begin{prop}[Intrinsic scaling]\label{prop:scaling-domain}
\locmark{loc:scaling-domain}
For $v\in\mathbb Q_{>0}$, define
\[
\mathcal D_v=\{\Phi\in\mathcal C_{\mathcal A}:vE(\Phi)\subseteq\mathbb Z\},
\qquad
\hat s_v\Phi=\Phi_{f_\Phi}^{\,vE(\Phi)}\quad(\Phi\in\mathcal D_v).
\]
Then $\hat s_v$ is a bijection from $\mathcal D_v$ to
$\mathcal D_{1/v}$ with inverse $\hat s_{1/v}$.
If $v=p/q$ in lowest terms, $\Phi\in\mathcal D_v$ precisely when
$E(\Phi)\subseteq q\mathbb Z$.
For any representation $\Phi=\Phi_f^N$ such that $vN\subseteq\mathbb Z$,
\begin{equation}\label{eq:scaling-padded}
\hat s_v\Phi_f^N=\Phi_f^{vN}.
\end{equation}
Moreover, whenever both successive operations are defined,
\begin{equation}\label{eq:scaling-composition}
\hat s_u(\hat s_v\Phi)=\hat s_{uv}\Phi.
\end{equation}
\end{prop}
\begin{proof}
The essential pair is unique. Since $v>0$, the order of its sites is
unchanged, and the local rule remains irreducible after scaling. Therefore
$E(\hat s_v\Phi)=vE(\Phi)$ and
$f_{\hat s_v\Phi}=f_\Phi$. This proves the inverse and composition
statements. The divisibility condition follows from $\gcd(p,q)=1$.
For a padded representation, write $f=g\circ h$ with $g$ irreducible and
$h$ the projection onto the essential coordinates. If
$N=\{j_1<\cdots<j_n\}$, its essential subset is
$E=\{j_i:i\in H(h)\}$. Scaling preserves that ordered selection, so
$\Phi_f^{vN}=\Phi_g^{vE}$, proving
\eqref{eq:scaling-padded}.
\end{proof}
The composition formula is an equality of values on the domain of the
successive operations. It does not assert equality of their domains:
for example, $\hat s_2\hat s_{1/2}$ is defined only on
$\mathcal D_{1/2}$, whereas $\hat s_1$ is defined everywhere.
An integer scaling $\hat s_p$ is defined on all global maps, but for
$p>1$ its image is the proper subset $\mathcal D_{1/p}$.
Nullary rules cause no exception: their essential support is empty and
every scaling fixes them.

\begin{prop}[Product dynamics]\label{prop:product}
\locmark{loc:product}
Put $X=\mathcal A^{\mathbb Z}$ and let $p\geq1$ be an integer.
For $0\leq r<p$, define $\delta_r:X\to X$ by
$(\delta_rx)_i=x_{pi+r}$, and define
\[
s_p:X\longrightarrow X^p,\qquad
s_p(x)=(\delta_0x,\ldots,\delta_{p-1}x).
\]
Then $s_p$ is a homeomorphism and, for every CA global map $\Phi$,
\begin{equation}\label{eq:product-conjugacy}
s_p\circ\hat s_p\Phi=\Phi^{\times p}\circ s_p,
\qquad
\hat s_p\Phi=s_p^{-1}\circ\Phi^{\times p}\circ s_p,
\end{equation}
where $\Phi^{\times p}(x^{(0)},\ldots,x^{(p-1)})
=(\Phi x^{(0)},\ldots,\Phi x^{(p-1)})$.
Each $\delta_r$ is a continuous surjective factor map satisfying
\begin{equation}\label{eq:decimation-factor}
\delta_r\circ\hat s_p\Phi=\Phi\circ\delta_r.
\end{equation}
\end{prop}
\begin{proof}
Every integer has a unique form $pi+r$ with $0\leq r<p$.
Thus the inverse of $s_p$ interleaves the $p$ configurations by
$s_p^{-1}(x^{(0)},\ldots,x^{(p-1)})_{pi+r}=x_i^{(r)}$.
Both maps are continuous in the product topology. If
$\Phi=\Phi_f^E$ and $E=\{j_1<\cdots<j_m\}$, then
\[
(\delta_r\hat s_p\Phi(x))_i
=f(x_{pi+r+pj_1}\cdots x_{pi+r+pj_m})
=(\Phi(\delta_rx))_i.
\]
This proves both identities, including the nullary case. Surjectivity of
$\delta_r$ follows by filling the other residue classes arbitrarily.
\end{proof}
Equation~\eqref{eq:product-conjugacy} relates the scaled CA to
$p$ independent copies of the original CA. Conjugacy to one copy
need not hold. For example,
$\hat s_2\sigma=\sigma^2$, but $\sigma$ has $k$ fixed points and
$\sigma^2$ has $k^2$, so they are not topologically conjugate for $k\geq2$.
The joint map $s_p$ retains every coordinate. In contrast, each individual
$\delta_r$ loses the other residue classes when $p>1$.

\paragraph{Similarity of supports and global maps.}
Two finite integer sets are called similar, $N\sim M$, if
\[
M=b+\epsilon vN
\quad\text{for some }b\in\mathbb Q,\quad
v\in\mathbb Q_{>0},\quad\epsilon\in\{1,-1\}.
\]
Both $N$ and $M$ are required to be subsets of $\mathbb Z$.
For global maps, let $\Phi\sim\Psi$ mean that there are
$\Phi'\cong\Phi$, $\Psi'\cong\Psi$ and $v>0$ rational such that
$\Phi'\in\mathcal D_v$ and $\Psi'=\hat s_v\Phi'$.
This agrees with allowing finite sequences of congruences and admissible
scalings, as the following characterization proves.

\begin{prop}[Similarity as an equivalence relation]
\label{prop:similarity}\locmark{loc:similarity}
Let $(N,f)$ and $(M,g)$ be the essential pairs of $\Phi$ and $\Psi$.
Then $\Phi\sim\Psi$ if and only if there exist
$b\in\mathbb Q$, $v\in\mathbb Q_{>0}$,
$\epsilon\in\{1,-1\}$, and $\nu\in S_{\mathcal A}$ such that
\begin{equation}\label{eq:affine-similarity}
M=b+\epsilon vN,\qquad
g=\hat\nu\hat r^{(1-\epsilon)/2}f.
\end{equation}
Consequently $\sim$ is reflexive, symmetric, and transitive on
$\mathcal C_{\mathcal A}$, independently of any local representation.
\end{prop}
\begin{proof}
Every congruence has the indicated form with $v=1$ and integer
translation, and every admissible scaling has the form with $b=0$,
$\epsilon=1$, and $\nu=1$.
Their compositions have the form \eqref{eq:affine-similarity}; hence that
form is necessary for both the stated relation and its finite-sequence
version.

Conversely, suppose \eqref{eq:affine-similarity} holds and $N\ne\emptyset$.
Reflect and relabel $(N,f)$ to obtain $(K,g)$, where $K=\epsilon N$.
Translate by $-\min K$, apply $\hat s_v$, and then translate by $\min M$.
All operations are admissible, since
\[
v(K-\min K)=M-\min M\subseteq\mathbb Z.
\]
They give $(M,g)$ and use only one scaling between two congruences.
This also covers singleton supports. If $N=\emptyset$, then $M=\emptyset$
and state relabeling alone gives the required relation.

The identity parameters establish reflexivity. Inverting the affine map
and the rule transformation establishes symmetry. For transitivity,
composing the affine maps with parameters $(b,v,\epsilon)$ and
$(b',v',\epsilon')$ gives parameters
\[
(b'+\epsilon'v'b,\ v'v,\ \epsilon'\epsilon).
\]
State permutations commute with reversal, so the corresponding rule
transformations compose as well. The characterization then gives
transitivity. Representation independence follows from the uniqueness
of the essential pairs.
\end{proof}

For example, the ECA rule $W^3_{90}$ has essential pair
$(\{-1,1\},W^2_6)$, where $W^2_6(a_0a_1)=a_0\oplus a_1$.
Its support is not an integer multiple $pM$ with $p>1$, but a translation
makes the scaling apparent:
\[
\Phi^{[-1,1]}_{W^3_{90}}
=\sigma^{-1}\hat s_2\Phi^{\{0,1\}}_{W^2_6}
\sim\Phi^{\{0,1\}}_{W^2_6}
=\Phi^{[-1,1]}_{W^3_{102}}.
\]

\begin{prop}\label{prop:global-scaling}\locmark{loc:similarity-count}
For $N=[1,n]$, the number of similarity classes is
\begin{equation}\label{eq:similarity-count}
\begin{split}
|G_2(N)/{\sim}|={}&
\sum_{[M]\in\mathcal P(N)/{\sim}}
 |\overline L_2(|M|)/\mathcal S_2|\\
&+\sum_{\substack{[M]\in\mathcal P(N)/{\sim}\\M\notin\Sym}}
 \left(|\overline t_{|M|}^{\langle1\rangle}|
       +|\overline t_{|M|}^{\langle c\rangle}|\right).
\end{split}
\end{equation}
\end{prop}
\begin{proof}
Similarity preserves cardinality and reflection symmetry of finite
supports. For a support with at least two sites, an affine similarity
mapping it onto itself has scale 1, by comparison of diameters. Its
orientation-preserving part is the identity, and an
orientation-reversing self-map exists precisely when it is
reflection-symmetrical. Proposition~\ref{prop:similarity} therefore
gives exactly the same local-orbit splitting as in
Proposition~\ref{prop:global-shift}, now summed over similarity classes
of supports. Empty and singleton supports have trivial local reversal
and are already included in the first sum.
\end{proof}
The counts for $n=0,1,2$ remain $1,3,7$. For $n=3,4$, all two-site
supports become similar, and these are the only additional support
identifications. Therefore
\begin{align}
|G_2([1,3])/{\sim}|&=85-4=81,\label{eq:classes81}\\
|G_2([1,4])/{\sim}|&=16\,656-2\cdot4=16\,648.
\end{align}
Exhaustive enumeration confirms both values. The recurrence computation
for $n=5$ gives $1\,074\,020\,458$ similarity classes; it is not an
exhaustive enumeration of all five-input rules.

\begin{prop}\label{commutators}\locmark{loc:commutators}
Let $v=p/q>0$ with positive integers $p,q$. On $\mathcal D_v$,
\begin{equation}\label{eq:scaling-commutators}
\hat\nu\hat s_v=\hat s_v\hat\nu,\qquad
\hat r\hat s_v=\hat s_v\hat r,\qquad
\sigma^p\hat s_v=\hat s_v\sigma^q.
\end{equation}
In each equality the two sides have the same domain.
\end{prop}
\begin{proof}
State relabeling leaves the essential support unchanged, and reflection
replaces $E$ by $-E$. These operations preserve $\mathcal D_v$ and
commute with multiplication of the support by $v$.
For the last equality,
\[
p+vE=v(q+E),\qquad
v(q+E)\subseteq\mathbb Z\ \Longleftrightarrow\ vE\subseteq\mathbb Z.
\]
The two essential pairs are thus identical on exactly $\mathcal D_v$.
This proof does not reorder partially defined inverse scalings and does
not require the displayed fraction $p/q$ to be reduced.
\end{proof}

For a contiguous essential support $N=[1,n]$, a similarity from $N$ to
itself induces only the usual state-permutation and reversal action on
the local rule. Hence
\[
|\overline L_{\mathcal A}(n)/\mathcal S_{\mathcal A}|
=|\overline G_{\mathcal A}(N)/{\cong}|
=|\overline G_{\mathcal A}(N)/{\sim}|.
\]
For arbitrary (possibly reducible) rules on that interval,
\[
|\overline L_{\mathcal A}(n)/\mathcal S_{\mathcal A}|
\leq|G_{\mathcal A}(N)/{\sim}|
\leq|G_{\mathcal A}(N)/{\cong}|
\leq|L_{\mathcal A}(n)/\mathcal S_{\mathcal A}|.
\]
\end{revision}

\subsection{Elementary Cellular Automata}
Recall that elementary CAs are precisely those binary CAs with the neighbourhood $N = [-1,1]$,
and that their global rule space is denoted by $G_2(N)$.
We set $\mathcal{E} = G_2(N)$.
In this subsection we examine in more detail the partitions of $\mathcal{E}$ that arise from various equivalence relations.
It is well-known that under reflection and permutation,  \rev{the set $\mathcal{E}$ decomposes} into 88 equivalence classes,
in symbols, $|\mathcal{E} / \mathcal{S}_2| = 88$.
All equivalence relations discussed here will be coarser than the partition $\mathcal{E} / \mathcal{S}_2$.
In general, if $R$ and $S$ are equivalence relations on a set $A$, and $R$ is coarser than $S$,
then each equivalence class $C$ of $A / R$ is either itself an equivalence class of $A / S$ or the union of
two or more equivalence classes from $A / S$.
We specify the members of equivalence classes using their Wolfram codes.
Whenever a number $c$ appears in one of the equivalence classes below, it is to be interpreted as the global map $\Phi^{N}_{W^3_c}$.
Every equivalence relation considered in what follows is obtained by merging two or more of the following equivalence classes of
$\mathcal{E} / \mathcal{S}_2$ into larger ones.
\begin{align*}
& A_1 = \{3,17,63,119\}, & & A_2=\{5,95\}, & & A_3 = \{10,80,175,245\}, \\
& A_4 = \{12,68,207,221\}, & & A_5=\{15,85\}, & & A_6=\{34,48,187,243\}, \\
& A_7 = \{51\}, & & A_8=\{60,102,153,195\}, & & A_9 = \{90, 165\}, \\
& A_{10} = \{136, 192, 238, 252\}, & & A_{11}=\{160,250\}, & & A_{12}=\{170,240\}, \\
& A_{13}=\{204\}, & &  A_{14} = \{77\},  & & A_{15} = \{232 \}, \\
& A_{16} = \{ 23 \}, & & A_{17} = \{178\}, & & A_{18}=\{105\}, \\
& A_{19} = \{150\}. & & & &
\end{align*}
The first further equivalence relation that we will consider is the congruence relation $\cong$ that
adds the shift operator to permutation and reflection.
By Eq.~(\ref{eq:classes85}) this equivalence relation induces 85 equivalence classes.
The equivalence classes of $\mathcal{E} / \mathcal{S}_2$ that are merged to obtain the coarser equivalence classes of
$\mathcal{E} / {\cong}$ are the following, see also~\cite{RUIVO2018280}.
\begin{align*}
& B_1 = A_4 \cup A_6, & & B_2 = A_5 \cup A_7, & & B_3 = A_{12} \cup A_{13}.
\end{align*}
Next, we consider the equivalence relation $\sim$ on $\mathcal{E}$, which adds scaling to the set of
symmetry operations.
By Eq.~(\ref{eq:classes81}) the number of equivalence classes decreases to 81.
The equivalence classes of $\mathcal{E} / {\sim}$ that do not coincide with the equivalence classes of
$\mathcal{E} / \mathcal{S}_2$ are the following.
\begin{align*}
& C_1 = A_3 \cup A_4 \cup A_6, & & C_2=A_5 \cup A_7, & & C_3=A_8 \cup A_9,  \\
& C_4 = A_{10} \cup A_{11}, & & C_5 = A_1 \cup A_2, & & C_6 = A_{12} \cup A_{13}.
\end{align*}
Epperlein \cite{epperlein2015classification, epperlein2017}
classified the elementary CAs up to topological conjugacy.
Permutation and reflection are general topological conjugacies and yield
the well-known 88 equivalence classes.
Epperlein demonstrated that, in addition, there exist further topological conjugacies between
specific pairs of elementary CAs, and that apart from these there are no others.
The pairs
\begin{align*}
& ([15], [170]), & & ([23], [178]), & & ([77],[232])
\end{align*}
of equivalence classes of $\mathcal{E} / \mathcal{S}_2$ \rev{are conjugate
via}
\[
\vartheta: \{0,1\}^\mathbb{Z} \to \{0,1\}^\mathbb{Z},  \ \
\vartheta(x)_i =
\begin{cases}
cx_i & \text{if $i$ is odd,} \\
x_i & \text{if $i$ is even.}
\end{cases}
\]
\rev{For the first pair, $\vartheta$ sends rule 15 to rule 240,
which belongs to $[170]$; reflection sends rule 240 to rule 170.
In particular, a conjugacy between these two named representatives
is $r\circ\vartheta$.
Rules 90, 105, and 150 are bipermutive radius-one CAs: fixing the
other two inputs makes the local rule a permutation of either outer
input. The bipermutive conjugacy theorem identifies each with the
one-sided full shift on four symbols, so they are conjugate to one
another \cite{Kurka2003TopologicalAS,epperlein2015classification}.}
Let us denote the equivalence classes of the \rev{elementary CAs} that are obtained by taking
all topological conjugacies into account by $\mathcal{E} / \mathrm{TC}$.
Consequently, the equivalence classes of $\mathcal{E} / \mathrm{TC}$ that differ from equivalence classes of $\mathcal{E}/ \mathcal{S}_2$
are the following.
\begin{align*}
& D_1 = A_5 \cup A_{12}, & & D_2 = A_{14} \cup A_{15}, & & D_3 = A_{16} \cup A_{17},
& & D_4 = A_{9} \cup A_{18} \cup A_{19}.
\end{align*}
The partition $\mathcal{E}/\mathrm{TC}$ thus contains 83 equivalence classes.
\begin{revision}
Finally, put
\locmark{loc:join}
\[
\Omega=(\sim\!\upharpoonright_{\mathcal E})\vee
       (\mathrm{TC}\!\upharpoonright_{\mathcal E}),
\]
the join of the two equivalence relations on the finite set $\mathcal E$.
Explicitly, $\Phi\,\Omega\,\Psi$ means that there is a finite sequence
$\Phi=\Phi_0,\Phi_1,\ldots,\Phi_\ell=\Psi$ of maps in $\mathcal E$
such that each successive pair is either similar or topologically
conjugate. For example, rule 15 is topologically conjugate to rule 170,
and rule 170 is similar to rule 204, so rules 15 and 204 lie in the
same $\Omega$-class. All maps in the sequence must be elementary CAs;
allowing intermediate maps outside $\mathcal E$ could in principle
produce a different relation. We make no assertion about that
ambient join. Let $\mathcal E/\Omega$ denote the resulting partition.
\end{revision}

Then the equivalence classes of $\mathcal{E}/ \Omega$ that are distinct from those of
$\mathcal{E} / \mathcal{S}_2$ are the following.
\begin{align*}
& E_1 = A_3 \cup A_4 \cup A_6, & & E_2=A_5 \cup A_7 \cup A_{12} \cup A_{13}, & & E_3=A_8 \cup A_9 \cup A_{18} \cup A_{19}, \\
& E_4 = A_{10} \cup A_{11}, & & E_5 = A_1 \cup A_2, & & E_6 = A_{14} \cup A_{15}, \\
& E_7 = A_{16} \cup A_{17}. & & & &
\end{align*}
We see that 19 equivalence classes of $\mathcal{E} / \mathcal{S}_2$ are merged into
7 equivalence classes of $\mathcal{E} / \Omega$.
The other equivalence classes of $\mathcal{E} / \Omega$ are also equivalence classes of $\mathcal{E} / \mathcal{S}_2$.
\rev{Consequently, $\mathcal{E} / \Omega$ contains $88 - 19 + 7 = 76$ equivalence classes.
This calculation uses Epperlein's complete topological-conjugacy
classification as input; our enumeration checks the resulting finite
join, rather than independently proving that dynamical classification.}

\begin{revision}
\section{Dynamical Meaning and Scope}
\label{sec:extended}\locmark{loc:dynamics}
Our operators act on the space $\mathcal C_{\mathcal A}$ of global maps.
They should be distinguished from maps acting on the configuration space
$X=\mathcal A^{\mathbb Z}$.
State relabeling $\nu:X\to X$ and reflection $r:X\to X$ are
homeomorphisms; their induced operators on rule space are
$\hat\nu\Phi=\nu\Phi\nu^{-1}$ and $\hat r\Phi=r\Phi r$.
Both give topological conjugacies of the corresponding dynamical systems.
The configuration map $\nu$ is itself a CA. Reflection is not a CA,
because it satisfies $r\sigma=\sigma^{-1}r$ rather than commuting
with the shift.

Postcomposition by a shift is a different operation. Writing
$T_a(\Phi)=\sigma^a\Phi$, shift commutation gives
\begin{equation}\label{eq:moving-frame}
(T_a(\Phi))^t=\sigma^{at}\Phi^t\qquad(t\geq0).
\end{equation}
Thus it describes a change of moving frame. Although $T_a$ is a bijection
on rule space with inverse $T_{-a}$, it is generally not conjugation on
configuration space or an automorphism of the monoid of CA global maps.
Indeed, $T_a(\id)=\sigma^a$.
In particular, the identity CA and the shift are congruent but are not
topologically conjugate: the identity fixes every configuration, whereas
the shift fixes only the constant configurations.

For integer scaling, the exact dynamical statement is the product
conjugacy \eqref{eq:product-conjugacy}. It preserves all information
jointly through the interleaving homeomorphism $s_p$.
Selecting one residue class instead gives the non-injective factor map
\eqref{eq:decimation-factor} when $p>1$.
More generally, a continuous surjection $h:X\to Y$ is a factor map
from $(X,\Phi)$ to $(Y,\Psi)$ when $h\Phi=\Psi h$; without that
intertwining identity, a projection alone need not define closed dynamics
on $Y$. Coarse-graining often uses such non-injective projections.
Our similarity relation is a geometric equivalence of essential rules
and supports, not a claim of conjugacy to a single copy or of an
information-losing coarse-graining of configurations.

Topological conjugacy and similarity are incomparable even on ECAs.
Rules 204 and 170 are congruent and hence similar, but represent the
identity and the shift. Conversely, rules 15 and 170 are topologically
conjugate by alternating complementation followed by reflection, as
recalled above. Their
irreducible unary rules are respectively complementation and identity,
which no allowed state relabeling or geometric operation identifies.
The 76-class relation $\Omega$ is specifically the join of these two
relations on the ECA set.

We do not claim an exhaustive classification of all transformations that
send CAs to CAs. For example, when $\mathcal A=\mathbb Z/k\mathbb Z$,
adding a fixed symbol $a$ to each rule output is postcomposition by the
CA $D_a(x)_i=x_i+a$; it is distinct from conjugating by a state
permutation. Likewise, inverse time evolution for reversible CAs and
general block factors are outside the geometric equivalence studied here.

\section{Conclusions}
\label{sec:sum}
The standard essential-neighbourhood representation makes it possible to
combine local stabilizer counts with support geometry. The resulting
recurrences enumerate irreducible binary rules, while the global counting
formula records the extra orientation classes arising from asymmetric
supports. On three- and four-site intervals this gives 85 and 16\,656
congruence classes, with independent exhaustive verification.

Rational scaling is well-defined on the intrinsic domain
$\mathcal D_v$, and the affine characterization proves that the induced
similarity relation is an equivalence relation. It gives 81 and 16\,648
classes on the same two intervals. Scaling by an integer $p$ has the
precise dynamical interpretation of $p$ independent copies; conjugacy
to one copy need not hold. Combining similarity with the known topological
conjugacies within the finite ECA set gives 76 classes. These counts
describe the chosen geometric relations and their finite join; they do
not constitute a general dynamical classification of cellular automata.
\end{revision}

\textbf{Credit authorship contribution statement:}
Martin Schaller: Writing---original draft, Visualization, Software, Investigation, Formal analysis, Conceptualization. Karl Svozil: Methodology, Investigation, Conceptualization.

\textbf{Declaration of competing interest:}
The authors declare that they have no known competing financial interests or personal relationships that could have appeared to influence the work reported in this paper.

\rev{\textbf{Data and code availability:} The supplementary verification package contains two independent exhaustive enumeration programs and their outputs for binary arities up to four, recurrence checks through $n=5$, and the complete list of four-site DEC-class mergers. No empirical data were used.}

\textbf{Acknowledgements:}
This research was funded in whole or in part by the Austrian Science Fund (FWF) [Grant DOI:10.55776/PIN5424624]. The authors acknowledge TU Wien Bibliothek for financial support through its Open Access Funding Programme.


\begin{revision}
\appendix
\section{Reproducible Enumeration}
\label{app:verification}\locmark{loc:verification}
The supplementary program uses the Wolfram coding convention stated in
Section~\ref{sec:CA}. For each arity $0\leq n\leq4$, it visits every
integer truth table $0\leq F<2^{2^n}$ and performs the following steps.
\begin{enumerate}
\item A coordinate is essential precisely when some pair of input words
differing only at that coordinate have different outputs. Select all
essential coordinates and form the reduced truth table $f$ by setting
the unused inputs to zero. Independence makes this choice immaterial.
The support $E$ is recorded as a sorted subset of $[0,n-1]$.
\item If $E$ is nonempty, normalize it to $E_0=E-\min E$.
For similarity, and $|E|\geq2$, further divide by
$d=\gcd\{e:e\in E_0\}$.
The empty support stays empty and a singleton normalizes to $\{0\}$.
\item With normalized support $P$ and its reflection
$P^*=\max P-P$ (both empty when $P$ is empty), use the lexicographic
minimum of the four pairs
\[
(P,f),\quad(P,\hat cf),\quad(P^*,\hat rf),\quad(P^*,\widehat{rc}f)
\]
as the canonical key. Count the distinct keys.
\end{enumerate}
For congruence this is a complete invariant by
Proposition~\ref{prop:global-eq}. For similarity,
Proposition~\ref{prop:similarity} applies: if $P$ and $Q$ are integer
sets of minimum zero with gcd 1 and $Q=vP$ for positive rational $v$,
then the denominator of $v$ divides every element of $P$, and applying
the inverse argument forces $v=1$. Thus gcd normalization removes
exactly the positive rational scale, with reflection still coupled to
the rule. This proves the correctness of the canonical-key method.

\begin{table}[htbp]
\begin{revision}
\centering
\caption{Independent exhaustive counts. The local column uses
reflection and state complementation on all $n$-input rules.}
\label{tab:verification}
\begin{tabular}{r|r|r|r|r}
$n$ & Truth tables & Local classes & Congruence & Similarity\\\hline
0 & 2 & 1 & 1 & 1\\
1 & 4 & 3 & 3 & 3\\
2 & 16 & 7 & 7 & 7\\
3 & 256 & 88 & 85 & 81\\
4 & 65\,536 & 16\,704 & 16\,656 & 16\,648
\end{tabular}
\par
\end{revision}
\end{table}
The same enumeration records irreducible orbits and their stabilizers
and agrees with every entry of Table~\ref{tab:s2_fs} through $n=4$.
The $n=5$ entries use algebraic fixed-point counts and the recurrences,
not exhaustive enumeration of the $2^{32}$ truth tables.

\begin{samepage}
For reproducibility, the unrestricted fixed-point counts used by the
program are
\[
\begin{aligned}
F_1&=2^{2^n}, &
F_r&=2^{(2^n+2^{\lceil n/2\rceil})/2},\\
F_c&=\begin{cases}0&n=0,\\2^{2^{n-1}}&n\geq1,\end{cases} &
F_{rc}&=\begin{cases}F_c&n\text{ odd},\\0&n\text{ even},\end{cases}\\
F_{\mathcal S_2}&=
\begin{cases}2^{(2^n+2^{\lceil n/2\rceil})/4}&n\text{ odd},\\
0&n\text{ even}.
\end{cases}
\end{aligned}
\]
\end{samepage}
These are the unrestricted counts underlying
\cite{SCHALLER2025100298}. The type counts follow by subtracting the
rules with full stabilizer and dividing by the orbit size:
\[
|t_n^{\mathcal S_2}|=F_{\mathcal S_2},\qquad
|t_n^{\langle a\rangle}|=(F_a-F_{\mathcal S_2})/2
\quad(a=r,c,rc),
\]
\[
|t_n^{\langle1\rangle}|
=(F_1-F_r-F_c-F_{rc}+2F_{\mathcal S_2})/4.
\]
Applying Proposition~\ref{prop:recurrences} supplies an independent
algebraic check on the direct stabilizer enumeration.
The program also checks \eqref{eq:projection-stabilizer} for every
truth table through arity four and verifies the explicit alternating
complementation conjugacies on both site parities. The finite ECA join
uses Epperlein's conjugacy identifications as input.

\paragraph{Direct shift-graph check of Ref.\ \cite{RUIVO2018280}.}
The second program visits padded $n$-input truth tables and constructs
three partitions with a disjoint-set algorithm: shifts alone (SEC),
reversal and state complementation (DEC), and all these operations (DECx).
For a positive displacement $t<n$, a shift is admissible precisely when
the last $t$ arguments are unused; for a negative displacement $-t$, the
first $t$ arguments must be unused. The shifted truth table is evaluated
by moving the input coordinates and filling out-of-range positions with
zero. The independence check guarantees that the filling is irrelevant.
The displacement zero adds no edge. This tests every admissible shift
in Definition~6 of \cite{RUIVO2018280}; their Definition~7 supplies the
other generators. The resulting counts are
\[
\begin{array}{c|rrr}
n&\mathrm{SEC}&\mathrm{DEC}&\mathrm{DECx}\\\hline
2&14&7&7\\
3&242&88&85\\
4&65\,282&16\,704&16\,656.
\end{array}
\]
For $n=4$ the output records, for every merger, the minimum Wolfram code
of each constituent DEC class. It contains 46 pairs and one triple,
accounting for all 48 reductions from DEC to DECx.
\end{revision}

\begin{revision}
\begin{thebibliography}{10}

\bibitem{10.1007/978-3-540-73001-9_1}
Luigi Acerbi, Alberto Dennunzio, and Enrico Formenti.
\newblock Shifting and lifting of cellular automata.
\newblock In S.~Barry Cooper, Benedikt L{\"o}we, and Andrea Sorbi, editors,
  {\em Computation and Logic in the Real World}, pages 1--10, Berlin,
  Heidelberg, 2007. Springer Berlin Heidelberg.
\newblock \doi{10.1007/978-3-540-73001-9_1}.

\bibitem{armstrong1997groups}
M.A. Armstrong.
\newblock {\em Groups and Symmetry}.
\newblock Undergraduate Texts in Mathematics. Springer New York, 1997.

\bibitem{Castillo-Ramirez03072023}
A.~Castillo-Ramirez, M.~Sanchez-Alvarez, A.~Vazquez-Aceves, and
  A.~Zaldivar-Corichi.
\newblock A generalization of cellular automata over groups.
\newblock {\em Communications in Algebra}, 51(7):3114--3123, 2023.
\newblock \doi{10.1080/00927872.2023.2177663}.

\bibitem{CASTILLORAMIREZ2020104533}
Alonso Castillo-Ramirez and Maximilien Gadouleau.
\newblock Elementary, finite and linear {vN}-regular cellular automata.
\newblock {\em Information and Computation}, 274:104533, 2020.
\newblock AUTOMATA 2017. \doi{10.1016/j.ic.2020.104533}.

\bibitem{CastilloVeliz2024}
Alonso Castillo-Ramirez and Eduardo Veliz-Quintero.
\newblock On the minimal memory set of cellular automata.
\newblock \url{https://arxiv.org/abs/2404.06394v2}, 2024.
\newblock arXiv:2404.06394, version 2.

\bibitem{CATTANEO19971593}
G.~Cattaneo, E.~Formenti, L.~Margara, and G.~Mauri.
\newblock Transformations of the one-dimensional cellular automata rule space.
\newblock {\em Parallel Computing}, 23(11):1593--1611, 1997.
\newblock Cellular automata. \doi{10.1016/S0167-8191(97)00076-8}.

\bibitem{CeccheriniCoornaert2010}
Tullio Ceccherini-Silberstein and Michel Coornaert.
\newblock {\em Cellular Automata and Groups}.
\newblock Springer Monographs in Mathematics. Springer, Berlin, Heidelberg,
  2010.
\newblock \doi{10.1007/978-3-642-14034-1}.

\bibitem{epperlein2015classification}
Jeremias Epperlein.
\newblock Classification of elementary cellular automata up to topological
  conjugacy.
\newblock In Jarkko Kari, editor, {\em Cellular Automata and Discrete Complex
  Systems}, volume 9099 of {\em Lecture Notes in Computer Science}, pages
  99--112. Springer, 2015.
\newblock \doi{10.1007/978-3-662-47221-7_8}.

\bibitem{epperlein2017}
Jeremias Epperlein.
\newblock {\em Topological Conjugacies Between Cellular Automata}.
\newblock PhD thesis, Technische Universit{\"a}t Dresden, 2017.

\bibitem{GARCIAMORALES2013276}
V.~García-Morales.
\newblock Symmetry analysis of cellular automata.
\newblock {\em Physics Letters A}, 377(3):276--285, 2013.
\newblock \doi{10.1016/j.physleta.2012.11.052}.

\bibitem{Graham94}
Ronald~L. Graham, Donald~E. Knuth, and Oren Patashnik.
\newblock {\em Concrete Mathematics: A Foundation for Computer Science}.
\newblock Addison-Wesley Longman Publishing Co., Inc., USA, 2nd edition, 1994.

\bibitem{DBLP:journals/mst/Hedlund69}
G.~A. Hedlund.
\newblock Endomorphisms and automorphisms of the shift dynamical systems.
\newblock {\em Math. Syst. Theory}, 3(4):320--375, 1969.
\newblock \doi{10.1007/BF01691062}.

\bibitem{Kurka2003TopologicalAS}
Petr K{\r u}rka.
\newblock {\em Topological and Symbolic Dynamics}, volume~11 of {\em Cours
  Sp{\'e}cialis{\'e}s}.
\newblock Soci{\'e}t{\'e} Math{\'e}matique de France, Paris, 2003.

\bibitem{Li1990TheSO}
Wentian Li and Norman~H. Packard.
\newblock The structure of the elementary cellular automata rule space.
\newblock {\em Complex Systems}, 4:281--297, 1990.

\bibitem{Lind_Marcus_1995}
Douglas Lind and Brian Marcus.
\newblock {\em An Introduction to Symbolic Dynamics and Coding}.
\newblock Cambridge University Press, Cambridge, 1995.

\bibitem{rotman2012introduction}
J.~J. Rotman.
\newblock {\em An Introduction to the Theory of Groups}.
\newblock Graduate Texts in Mathematics. Springer, New York, NY, 2012.
\newblock \doi{10.1007/978-1-4612-4176-8}.

\bibitem{RUIVO2018280}
Eurico L.~P. Ruivo, Pedro P.~B. {de Oliveira}, Fabiola Lobos, and Eric Goles.
\newblock Shift-equivalence of k-ary, one-dimensional cellular automata rules.
\newblock {\em Communications in Nonlinear Science and Numerical Simulation},
  63:280--291, 2018.
\newblock \doi{10.1016/j.cnsns.2018.03.017}.

\bibitem{SCHALLER2025100298}
Martin Schaller and Karl Svozil.
\newblock Orbits of one-dimensional cellular automata induced by symmetry
  transformations.
\newblock {\em Physics Open}, 24:100298, 2025.
\newblock \doi{10.1016/j.physo.2025.100298}.

\bibitem{VISPOEL2022133074}
Milan Vispoel, Aisling~J. Daly, and Jan~M. Baetens.
\newblock Progress, gaps and obstacles in the classification of cellular
  automata.
\newblock {\em Physica D Nonlinear Phenomena}, 432:133074, 2022.
\newblock \doi{10.1016/j.physd.2021.133074}.

\bibitem{wolfram-1986}
S.~Wolfram.
\newblock {\em Theory and Applications of Cellular Automata: Including Selected
  Papers, 1983-1986}.
\newblock Advanced series on complex systems. World Scientific, 1986.

\end{thebibliography}
\end{revision}


\end{document}
